%% ============================================================================
%% Chapter 2: Consuming APIs in Python: From urllib to httpx
%% Mastering APIs: From Zero to Production Guru
%% ============================================================================

\chapter{Consuming APIs in Python: From urllib to httpx}
\label{ch:consuming_python}

%% ---------------------------------------------------------------------------
%% 1. Executive Summary
%% ---------------------------------------------------------------------------

Consuming an HTTP API looks deceptively simple: build a URL, send a request,
read some JSON. That mental model survives exactly until the first production
incident. A consumer that omits timeouts will, under a slow dependency, exhaust
its worker pool and take down a healthy service. A consumer that retries
blindly will multiply load on an already struggling server and convert a brown
outage into a full one. A consumer that parses a response body before checking
the status code will turn a perfectly meaningful \texttt{503} into a confusing
\texttt{JSONDecodeError} three stack frames away from the real cause. This
chapter treats API consumption as the engineering discipline it actually is.

We trace the evolution of Python HTTP clients --- from the standard library's
\texttt{urllib}, through the ergonomic revolution of \texttt{requests}
\cite{requestsdocs}, to the modern \texttt{httpx} library with its unified
sync/async API and HTTP/2 support \cite{httpxdocs} --- and extract the design
lessons each generation teaches. We then build, piece by piece, a
production-grade client: explicit connect/read/write/pool timeouts, bounded
connection pools with keep-alive, retry logic with exponential backoff and full
jitter (with the underlying mathematics derived, not hand-waved), honoring of
\texttt{Retry-After} and \texttt{RateLimit} headers per RFC~9110
\cite{rfc9110}, typed pagination iterators, streaming downloaders with
cancellation, and an error taxonomy that maps transport failures, status
failures, and contract failures to distinct, actionable exceptions. Every
listing is fully typed, runs on Python~3.11+, and is validated \emph{offline}
against a local stub server or \texttt{httpx.MockTransport}: no network, no
cloud accounts, no API keys, no flakiness.

The capstone of the chapter is a resilient SDK wrapper around the fictional
\emph{Northwind Robotics Parts API}: typed \texttt{pydantic} response models
\cite{pydanticdocs}, telemetry hooks, transport injection for hermetic tests,
and a complete fault-matrix test suite simulating \texttt{500}s, \texttt{429}s
with \texttt{Retry-After}, and slow responses. The resilience philosophy
throughout follows Nygard's \emph{Release It!} \cite{nygard2018releaseit}:
every integration point is a betrayal waiting to happen, so engineer the
client to fail fast, fail visibly, and recover gracefully.

%% ---------------------------------------------------------------------------
%% 2. Learning Objectives & Prerequisites
%% ---------------------------------------------------------------------------

\section{Learning Objectives \& Prerequisites}

\subsection{Learning Objectives}

After completing this chapter, you will be able to:

\begin{enumerate}
  \item \textbf{Select} the appropriate Python HTTP client library
        (\texttt{urllib}, \texttt{requests}, \texttt{httpx}) for a given
        workload, justifying the choice in terms of sync vs.\ async support,
        HTTP/2 capability, and dependency footprint.
  \item \textbf{Design} a robust client anatomy: base-URL joining, header and
        authentication injection, event hooks, and transport abstraction.
  \item \textbf{Specify} a complete timeout policy (connect, read, write,
        pool) and explain why each dimension exists and what failure it
        bounds.
  \item \textbf{Size} connection pools correctly, reasoning about keep-alive,
        DNS caching, TLS session reuse, and the cost of connection churn.
  \item \textbf{Derive and implement} exponential backoff with full jitter,
        compute expected-delay bounds, enforce retry budgets, honor
        \texttt{Retry-After}, and gate retries behind idempotency analysis.
  \item \textbf{Implement} client-side rate limiting (token bucket) that
        composes cleanly with server-driven \texttt{429} signals.
  \item \textbf{Consume} offset, cursor, and link-header paginated endpoints
        as lazy typed iterators, and stream large payloads with progress
        reporting and cancellation.
  \item \textbf{Classify} consumer-side errors into transport, status, and
        contract categories and map each to an actionable exception hierarchy.
  \item \textbf{Build and test} a typed SDK wrapper with telemetry hooks and
        hermetic offline tests via \texttt{httpx.MockTransport}.
\end{enumerate}

\subsection{Prerequisites}

This chapter assumes mastery of:

\begin{itemize}
  \item \textbf{Chapter 0:} the curriculum's tooling baseline --- Python 3.11+,
        a virtual environment, \texttt{pip}, and PowerShell on Windows.
  \item \textbf{Chapter 1:} the HTTP wire protocol --- methods, status-code
        classes, headers, message bodies, content negotiation, and the
        semantics of idempotent vs.\ non-idempotent methods as standardized in
        RFC~9110 \cite{rfc9110}.
  \item Core Python: type hints, context managers (\texttt{with}), iterators
        and generators, and the \texttt{logging} module.
\end{itemize}

No third-party package knowledge is assumed; each library is introduced from
first principles. All code runs offline against the local stub server provided
in Section~\ref{sec:ch02-lab} or against \texttt{httpx.MockTransport}.

\begin{keyconceptbox}
\textbf{The consumer's contract.} A production API client is not ``code that
calls an endpoint.'' It is a managed boundary between your process and an
unreliable network dependency. Its obligations are: bound every wait
(timeouts), bound every resource (pools, concurrency, memory), bound every
retry (budgets, jitter, idempotency), and make every failure observable and
actionable. Everything in this chapter serves one of those four bounds.
\end{keyconceptbox}

%% ---------------------------------------------------------------------------
%% 3. Deep Technical Mechanics & Architectural Concepts
%% ---------------------------------------------------------------------------

\section{Deep Technical Mechanics \& Architectural Concepts}

\subsection{The Evolution of Python HTTP Clients}
\label{sec:ch02-evolution}

\subsubsection{urllib: the stdlib baseline}

Python ships with \texttt{urllib.request}, a functional but austere client.
Its virtues are real: zero third-party dependencies, availability on every
interpreter, and therefore suitability for constrained environments
(installers, bootstrap scripts, single-file tools, frozen binaries). Its sharp
edges are equally real:

\begin{itemize}
  \item \textbf{A single global timeout} passed per call --- no distinction
        between connect and read phases, and no pool-acquisition bound.
  \item \textbf{No session abstraction.} Every call re-resolves DNS,
        re-opens TCP, and re-performs the TLS handshake unless you hand-manage
        an \texttt{OpenerDirector}; keep-alive is not the default behavior of
        the high-level API.
  \item \textbf{Exceptions for non-2xx statuses.} \texttt{HTTPError} is
        raised for 4xx/5xx, which is defensible but surprises developers
        trained on \texttt{requests}, and the error doubles as a response
        object, a notorious source of confusion.
  \item \textbf{Manual everything:} JSON encoding/decoding, query-string
        construction, redirect policy, decompression negotiation, and proxy
        configuration are all the caller's problem.
\end{itemize}

The lesson \texttt{urllib} teaches is foundational: \emph{an HTTP client is a
state machine over sockets}, and if the library does not manage that state for
you, your application code will --- badly, and usually at 3~a.m.

\subsubsection{requests: the ergonomic revolution}

\texttt{requests} \cite{requestsdocs} re-founded Python HTTP consumption on
two ideas: a humane API (``HTTP for Humans'') and the \textbf{Session} object.
A \texttt{Session} persists cookies, default headers, and --- critically ---
a connection pool across requests, making keep-alive the default. Its
limitations only became pressing a decade later: it is synchronous only,
HTTP/1.1 only, built on \texttt{urllib3}, and its timeout model (a single
value or a \texttt{(connect, read)} pair) does not cover write or
pool-acquisition phases. For a synchronous, HTTP/1.1, request/response
workload, \texttt{requests} remains a correct and battle-tested choice; its
deficit is architectural headroom, not correctness.

\subsubsection{httpx: the modern synthesis}

\texttt{httpx} \cite{httpxdocs} keeps the \texttt{requests}-shaped API and
adds the features the modern era demands:

\begin{itemize}
  \item \textbf{Dual sync/async clients} (\texttt{httpx.Client} and
        \texttt{httpx.AsyncClient}) sharing one API surface --- the same
        instrumentation, retry, and pagination code paths serve threads and
        \texttt{asyncio} alike.
  \item \textbf{HTTP/2 support} (via the optional \texttt{h2} package),
        enabling multiplexed streams over a single connection.
  \item \textbf{A four-dimensional timeout model}: connect, read, write, and
        pool timeouts are independently configurable.
  \item \textbf{Transport injection}: \texttt{httpx.MockTransport} lets tests
        exercise the \emph{entire} client stack --- pooling, redirects,
        timeouts, hooks --- without a socket ever being opened.
  \item \textbf{Event hooks}: per-client middleware for logging, tracing, and
        telemetry.
\end{itemize}

\begin{table}[h]
  \centering
  \small
  \begin{tabularx}{\textwidth}{l c c c X}
    \toprule
    \textbf{Capability} & \textbf{urllib} & \textbf{requests} & \textbf{httpx} & \textbf{Notes} \\
    \midrule
    Stdlib (no deps)      & Yes & No  & No  & urllib wins constrained environments \\
    Session / pooling     & Manual & Yes & Yes & keep-alive default in both libraries \\
    Async support         & No  & No  & Yes & httpx: asyncio and trio \\
    HTTP/2                & No  & No  & Opt-in & requires \texttt{h2} extra \\
    Timeout dimensions    & 1   & 2   & 4   & httpx adds write + pool \\
    Mock transport        & No  & Adapter-based & First-class & MockTransport runs the full stack \\
    Type hints            & Partial & Partial & Full & httpx ships \texttt{py.typed} \\
    \bottomrule
  \end{tabularx}
  \caption{HTTP client support matrix for Python 3.11+ consumers.}
  \label{tab:ch02-client-matrix}
\end{table}

\textbf{Selection heuristic.} Use \texttt{urllib} only when third-party
dependencies are forbidden and the call count is tiny. Use \texttt{requests}
when the codebase is synchronous, HTTP/1.1 suffices, and the dependency
already exists. Use \texttt{httpx} for anything new: async workloads, HTTP/2
upstreams, or any system where hermetic testing matters --- which is to say,
any system. The remainder of this chapter standardizes on \texttt{httpx}.

\subsection{Anatomy of a Robust Client}
\label{sec:ch02-anatomy}

A well-built client is a pipeline, not a function call.
Figure~\ref{fig:ch02-lifecycle} shows the stages every request traverses in
the production client we build in Section~\ref{sec:ch02-code}.

\begin{figure}[h]
\centering
\begin{tikzpicture}[
    font=\small,
    stage/.style={rectangle, draw=darkblue, fill=blue!8, thick,
                  minimum width=2.05cm, minimum height=1.0cm, align=center},
    sidebox/.style={rectangle, draw=packtgray, fill=lightgray, thick,
                    minimum width=2.05cm, minimum height=0.9cm, align=center},
    flow/.style={->, thick, >=stealth},
    retryflow/.style={->, thick, >=stealth, dashed, draw=packtorange}
]
  % Main pipeline
  \node[stage] (build) {Build\\request};
  \node[stage, right=0.55cm of build] (auth) {Inject auth\\\& headers};
  \node[stage, right=0.55cm of auth] (pool) {Acquire pooled\\connection};
  \node[stage, right=0.55cm of pool] (tls) {TLS handshake\\(if cold)};
  \node[stage, right=0.55cm of tls] (send) {Send \&\\await response};
  \node[stage, right=0.55cm of send] (parse) {Status check\\\& typed parse};

  \draw[flow] (build) -- (auth);
  \draw[flow] (auth) -- (pool);
  \draw[flow] (pool) -- (tls);
  \draw[flow] (tls) -- (send);
  \draw[flow] (send) -- (parse);

  % Retry loop
  \node[sidebox, below=1.15cm of send] (retry) {Retry?\\budget / jitter /\\idempotency gate};
  \draw[retryflow] (send.south) -- node[right, align=left, xshift=2pt]
       {\footnotesize retryable failure} (retry.north);
  \draw[retryflow] (retry.west) -| node[pos=0.25, above, align=center]
       {\footnotesize backoff sleep} ($(pool.south)!0.5!(tls.south)$);

  % Timeout brace above
  \draw[decorate, decoration={brace, amplitude=6pt}, thick, draw=packtgray]
    ($(pool.north west)+(0,0.18)$) -- ($(send.north east)+(0,0.18)$)
    node[midway, above=8pt, align=center] {\footnotesize timeouts bound every phase};

  % Observability side box
  \node[sidebox, below=1.15cm of pool] (obs) {Telemetry hooks\\(logging, timing)};
  \draw[flow, draw=packtgray] (obs.north) -- (pool.south);
  \draw[flow, draw=packtgray] (obs.east) -| (parse.south);
\end{tikzpicture}
\caption{Client request lifecycle pipeline. Every phase between pool
acquisition and response receipt is bounded by an explicit timeout; failures
route through a retry gate that enforces budget, jitter, and idempotency;
telemetry hooks observe every stage.}
\label{fig:ch02-lifecycle}
\end{figure}

The pipeline's constituent responsibilities:

\begin{description}
  \item[Base URL handling.] Store the base URL once; join paths with a
        \texttt{httpx.URL} (or \texttt{urljoin} semantics) rather than string
        concatenation. Trailing-slash bugs (\texttt{api.example.com/v1} +
        \texttt{/parts}) are the single most common client configuration
        defect. \texttt{httpx.Client(base\_url=...)} implements the correct
        merge semantics.
  \item[Header management.] Set protocol headers (\texttt{Accept},
        \texttt{Accept-Encoding}, \texttt{User-Agent}) at client construction.
        A descriptive \texttt{User-Agent} (product, version, contact) is
        professional courtesy --- and often a provider requirement --- so
        your traffic is identifiable when you are accidentally the source of
        an incident.
  \item[Authentication injection.] Credentials live in exactly one place: an
        \texttt{auth} callable or a constructor-level header. Never sprinkle
        \texttt{headers=\{"Authorization": ...\}} across call sites, and
        never log the \texttt{Authorization} header (see
        Section~\ref{sec:ch02-pitfalls}).
  \item[Hooks / middleware.] \texttt{httpx} event hooks
        (\texttt{request}, \texttt{response}) are the seam for telemetry,
        trace-context propagation, and cross-cutting policy --- without
        polluting business logic.
\end{description}

\subsection{Timeout Taxonomy: Bounding Every Wait}
\label{sec:ch02-timeouts}

\begin{definition}[Timeout]
A \emph{timeout} is an upper bound on the wall-clock time a specific phase of
an operation may consume before the operation is aborted and the phase's
resources are reclaimed. A timeout is a correctness property, not a
performance tuning knob: it converts an unbounded wait (a hang) into a bounded
failure (an exception).
\end{definition}

A request without a timeout is a thread (or task) that can be held
\emph{forever}. Nygard documents the canonical failure mode
\cite{nygard2018releaseit}: a dependency degrades until responses slow to a
crawl; every worker in the calling service blocks awaiting sockets that will
never answer; the caller's pool exhausts; the caller stops serving its own
traffic; the outage propagates upstream. The dependency's partial failure
becomes the consumer's total failure --- a cascading failure enabled entirely
by the missing timeout. Hence the industry adage we elevate to law:

\begin{warningbox}
\textbf{There is no safe default of ``no timeout.''} \texttt{urllib}'s global
default is \emph{infinite}; \texttt{requests} defaults to \emph{infinite};
\texttt{httpx} defaults to 5~seconds and must be explicitly disabled to
remove. Any code review that finds a bare \texttt{requests.get(url)} or
\texttt{urlopen(url)} with no timeout has found a latent severity-1 incident.
``It has never hung'' is a statement about the past, not a property of the
system.
\end{warningbox}

\subsubsection{The four timeout dimensions}

\texttt{httpx.Timeout} exposes the complete taxonomy
\cite{httpxdocs}:

\begin{description}
  \item[connect] --- time to establish the connection: DNS resolution, TCP
        handshake, and TLS negotiation. Bounds a dead or firewalled host.
        Typical: 1--5~s. Connects to a healthy host on a healthy network
        complete in tens of milliseconds; a 5~s connect timeout is already
        generous.
  \item[read] --- maximum gap while awaiting bytes from the peer: time to
        first byte \emph{and} inter-chunk gaps while streaming a body. Bounds
        a hung application behind a healthy load balancer --- the failure mode
        TCP cannot see, because the connection is fine and the peer simply
        never sends. This is the timeout whose absence causes the war story in
        Section~\ref{sec:ch02-warstory}. Typical: 5--30~s, aligned to the
        endpoint's documented p99.
  \item[write] --- maximum gap while sending the request body. Matters for
        uploads and large POST bodies; bounds a peer that stops reading.
        Typical: 5--30~s.
  \item[pool] --- time to wait for a connection from a saturated pool. Bounds
        queueing delay when concurrency exceeds pool size, turning silent
        queueing (head-of-line blocking that looks like a hang) into a loud,
        actionable \texttt{PoolTimeout}. Typical: 1--5~s.
\end{description}

\textbf{Per-call vs.\ client defaults.} Set a conservative default on the
client (e.g., connect 2~s, read 10~s, write 10~s, pool 2~s) and tighten or
loosen per call where the endpoint contract justifies it --- a report-generation
endpoint may legitimately need a 120~s read timeout, but that exception must
be explicit at the call site, not global. A default of ``whatever the library
ships'' is an abdication; a default of ``infinite'' is a defect.

\subsection{Connection Pooling, Keep-Alive, and the Cost of Churn}
\label{sec:ch02-pooling}

Establishing a fresh HTTPS connection costs: one DNS lookup (typically
cached, 1--50~ms when not), one TCP handshake (1 RTT), and one TLS~1.3
handshake (1 RTT more, plus certificate validation and CPU for asymmetric
cryptography). On a cross-country path at 60~ms RTT, a cold connection adds
$\sim$120--200~ms before the first request byte flows --- and the full
handshake's key-exchange cost is paid again for every request. HTTP keep-alive
amortizes all of it: the pool holds idle connections open and hands them to
subsequent requests, reducing per-request overhead to the request/response
exchange itself.

\texttt{httpx.Limits} exposes the two sizing levers
\cite{httpxdocs}:

\begin{description}
  \item[max\_connections] --- total connections the pool will hold. Size it to
        the peak concurrent requests your process will issue to \emph{all}
        hosts. Default: 100.
  \item[max\_keepalive\_connections] --- connections retained idle for reuse.
        Size it to your steady-state concurrency against this client.
        Default: 20. A value too small causes \emph{churn}: connections are
        closed on response completion and rebuilt on the next request,
        re-paying handshake costs under load.
\end{description}

Three subtleties separate a correctly tuned pool from a cargo-culted one:

\begin{enumerate}
  \item \textbf{DNS caching.} The OS resolver cache is outside Python's
        control; long-lived pools naturally avoid re-resolution because
        keep-alive connections skip DNS entirely. For short-lived
        connections, resolution happens per connection, and a flaky resolver
        becomes a latency multiplier.
  \item \textbf{TLS session reuse.} TLS~1.3 session tickets allow
        \emph{resumed} handshakes (0-RTT/1-RTT reduction) on new connections
        to the same host; libraries built on the stdlib \texttt{ssl} module
        benefit when a shared \texttt{SSLContext} is used. Keep-alive still
        dominates: the cheapest handshake is the one you never perform.
  \item \textbf{One client per base URL, long-lived.} Create the
        \texttt{Client} once (module scope or dependency-injected), reuse it
        for the process lifetime, and close it deterministically (context
        manager or shutdown hook). Creating a client per request discards the
        pool --- the single most common httpx misuse --- while never closing
        clients leaks sockets until the OS reclaims them.
\end{enumerate}

\subsection{Retry Engineering: Backoff, Jitter, Budgets, Idempotency}
\label{sec:ch02-retries}

Retries exist because a large class of failures is \emph{transient}: packet
loss, a restarting pod, a load balancer draining a node, a brief 502 from a
proxy mid-deploy. Retrying converts transient failure into success. Retrying
\emph{wrong} converts a partial outage into a total one --- the retry storm of
Section~\ref{sec:ch02-warstory}.

\subsubsection{Exponential backoff and full jitter: the math}

Let attempts be numbered $n = 1, 2, \dots, N$ after the initial call, with
base delay $b$ and cap $c$. Uncapped exponential backoff waits
\begin{equation}
  W_n = b \cdot 2^{\,n-1},
\end{equation}
so with $b = 1\,\text{s}$ the schedule is $1, 2, 4, 8, 16, \dots$ seconds.
Exponential spacing buys the server recovery time, but it has a fatal flaw
under correlated failure: when an outage ends, \emph{every} queued client
wakes on the same schedule and re-hammers the server in synchronized waves ---
the \emph{thundering herd}.

\textbf{Full jitter} (popularized by the AWS Architecture Blog, see
Section~\ref{sec:ch02-bibliography}) draws each delay uniformly from a range:
\begin{equation}
  W_n \;\sim\; \mathrm{Uniform}\bigl(0,\; \min(c,\; b \cdot 2^{\,n-1})\bigr).
\end{equation}
Because a uniform variable on $[0, m]$ has mean $m/2$, the expected per-attempt
delay while the cap is inactive is
\begin{equation}
  \mathbb{E}[W_n] \;=\; \frac{b \cdot 2^{\,n-1}}{2},
  \qquad
  \mathbb{E}\!\left[\sum_{n=1}^{N} W_n\right]
  \;=\; \frac{b}{2}\left(2^{N} - 1\right),
  \label{eq:ch02-jitter-total}
\end{equation}
exactly half the deterministic schedule, with the same exponential growth
profile. Jitter sacrifices no meaningful aggressiveness while decorrelating
clients: collision probability across $k$ concurrent retriers drops from
$\sim$1 (deterministic) to $\sim$1/(range width). A variant,
\emph{decorrelated jitter}, samples $W_n \sim \mathrm{Uniform}(b,\,
3 W_{n-1}) \wedge c$, coupling consecutive waits instead of attempt indices;
it spreads load comparably and converges faster after short blips. Both are
vastly superior to unjittered or ``equal jitter'' schedules; full jitter is
the simplest to reason about and is what we implement.

\begin{figure}[h]
\centering
\begin{tikzpicture}
\begin{axis}[
    width=0.86\textwidth, height=6.2cm,
    xlabel={Retry attempt $n$},
    ylabel={Delay (seconds)},
    xmin=0.5, xmax=6.5, ymin=0, ymax=33,
    xtick={1,2,3,4,5,6},
    legend style={at={(0.02,0.97)}, anchor=north west, font=\footnotesize},
    grid=both, grid style={gray!25},
]
  % Deterministic exponential: b=1 -> 1,2,4,8,16,32
  \addplot[mark=*, thick, color=darkblue] coordinates
    {(1,1) (2,2) (3,4) (4,8) (5,16) (6,32)};
  \addlegendentry{Deterministic: $W_n = 2^{n-1}$}
  % Full-jitter expectation: b/2 * 2^{n-1}
  \addplot[mark=square*, thick, color=packtorange] coordinates
    {(1,0.5) (2,1) (3,2) (4,4) (5,8) (6,16)};
  \addlegendentry{Full jitter expectation: $\mathbb{E}[W_n] = 2^{n-2}$}
  % One seeded full-jitter sample realization (seed 42, illustrative)
  \addplot[mark=triangle*, dashed, thick, color=packtgray] coordinates
    {(1,0.64) (2,0.28) (3,3.15) (4,3.88) (5,11.6) (6,2.79)};
  \addlegendentry{Full jitter, one seeded realization}
\end{axis}
\end{tikzpicture}
\caption{Backoff schedules for $b = 1\,\text{s}$. Deterministic backoff grows
exponentially but synchronizes clients; full jitter halves the expected delay
and, crucially, decorrelates concurrent retriers (one seeded realization
shown). Production policy caps both at $c$ (here 30~s).}
\label{fig:ch02-backoff-curve}
\end{figure}

Figure~\ref{fig:ch02-retry-timeline} shows the same policy as a timeline for a
single request's life.

\begin{figure}[h]
\centering
\begin{tikzpicture}[
    font=\small,
    attempt/.style={rectangle, draw=darkblue, fill=blue!10, thick,
                    minimum height=0.8cm, align=center},
    sleep/.style={rectangle, draw=packtorange, fill=orange!12, thick,
                  minimum height=0.8cm, align=center},
    outcome/.style={circle, draw=packtgray, fill=green!12, thick,
                    inner sep=1.5pt, align=center}
]
  \node[attempt, minimum width=1.5cm] (a0) at (0,0) {Attempt 0\\\footnotesize 500};
  \node[sleep,   minimum width=1.1cm, right=0.35cm of a0] (s1) {jit\\\footnotesize 0.4s};
  \node[attempt, minimum width=1.5cm, right=0.35cm of s1] (a1) {Attempt 1\\\footnotesize 500};
  \node[sleep,   minimum width=1.7cm, right=0.35cm of a1] (s2) {jitter\\\footnotesize 1.7s};
  \node[attempt, minimum width=1.5cm, right=0.35cm of s2] (a2) {Attempt 2\\\footnotesize 429};
  \node[sleep,   minimum width=2.3cm, right=0.35cm of a2] (s3) {Retry-After\\\footnotesize 5.0s};
  \node[attempt, minimum width=1.5cm, fill=green!15, right=0.35cm of s3] (a3) {Attempt 3\\\footnotesize 200 OK};
  \node[outcome, right=0.45cm of a3] (ok) {\checkmark};

  \foreach \x/\y in {a0/s1, s1/a1, a1/s2, s2/a2, a2/s3, s3/a3, a3/ok}
    \draw[->, thick, >=stealth] (\x) -- (\y);

  \draw[decorate, decoration={brace, mirror, amplitude=6pt}, thick, draw=packtgray]
    ($(a0.south west)+(0,-0.12)$) -- ($(ok.south east)+(0,-0.12)$)
    node[midway, below=8pt, align=center]
    {\footnotesize retry budget: total time \& attempts bounded; 429 defers to server-advertised delay};
\end{tikzpicture}
\caption{Retry timeline for one request. Transient \texttt{500}s consume
full-jitter backoff slots; a \texttt{429} overrides the schedule with the
server's \texttt{Retry-After}; the budget caps total attempts and elapsed
time.}
\label{fig:ch02-retry-timeline}
\end{figure}

\subsubsection{Retry budgets}

Per-request caps ($N \le 3$--$5$ attempts, a total-time ceiling) bound the
blast radius of one call. They do not bound the \emph{fleet}: if 30\% of all
calls are failing, a retry-everything policy multiplies offered load by up to
$N$ and finishes the server off. A \textbf{retry budget} rate-limits retries
themselves --- e.g., ``no more than 10\% of issued requests per minute may be
retries'' --- typically implemented as a shared token bucket. When the budget
is empty, fail fast and let the circuit breaker (Chapter~15) take over; the
two patterns compose: retries handle transient blips, the breaker handles
sustained failure, the budget prevents retries from feeding the fire.

\subsubsection{Retry-After and which statuses merit retry}

RFC~9110 \cite{rfc9110} defines \texttt{Retry-After} (on \texttt{429} and
\texttt{503}) as the server's explicit request for a delay, in seconds or
HTTP-date. A server that tells you when to come back is giving you better
information than any locally computed backoff; \textbf{always honor it}
(clamped to a sane ceiling) when present. Retry policy by status class:

\begin{itemize}
  \item \textbf{Retry:} \texttt{408}, \texttt{425}, \texttt{429}, and
        \texttt{5xx} (with skepticism for \texttt{501}/\texttt{505}), plus
        transport errors (connect failures, connection resets, read
        timeouts).
  \item \textbf{Never retry (as-is):} \texttt{400}, \texttt{401},
        \texttt{403}, \texttt{404}, \texttt{409}, \texttt{422} --- these are
        deterministic; the identical request will fail identically.
\end{itemize}

\subsubsection{The idempotency gate}

A retry replays a request. Replaying is safe only when replay cannot change
the outcome: per RFC~9110 \cite{rfc9110}, \texttt{GET}, \texttt{HEAD},
\texttt{PUT}, \texttt{DELETE}, and \texttt{OPTIONS} are idempotent by method
semantics. \texttt{POST} is not --- a retried \texttt{POST /orders} may create
two orders, and a read timeout is the treacherous case: the server may have
processed the request \emph{and then} the response was lost, so the client's
``failure'' is the server's ``success.'' The professional escape hatches, in
order of preference:

\begin{enumerate}
  \item \textbf{Idempotency keys:} the API accepts a client-generated
        \texttt{Idempotency-Key} header and deduplicates replays server-side
        (Chapter~10 covers the provider side). With a key, retrying a
        \texttt{POST} becomes safe.
  \item \textbf{Natural idempotency:} the operation is semantically
        repeatable (``set status to X''), verifiable from the endpoint
        contract.
  \item \textbf{No retry:} when neither holds, propagate the failure and let
        a human or a reconciliation job decide. A client that guesses wrong
        here causes data corruption, the worst failure class there is.
\end{enumerate}

\subsection{Rate-Limit-Aware Consumption}
\label{sec:ch02-ratelimit}

Server-side rate limiting (Chapter~14) protects the provider; client-side
throttling is the consumer's half of the contract. A well-behaved client:

\begin{enumerate}
  \item \textbf{Proactively paces itself} below the published limit with a
        token bucket: a bucket of capacity $B$ refills at rate $r$ tokens/s;
        each request consumes one token; an empty bucket blocks until a token
        accrues. Sustained throughput converges to $r$ with burst tolerance
        $B$.
  \item \textbf{Reacts to signals:} on \texttt{429}, read \texttt{Retry-After}
        and the \texttt{RateLimit-*} header family (limit, remaining, reset),
        back the bucket off, and retry only after the advertised window.
  \item \textbf{Never compensates for a 429 by raising parallelism} --- the
        queue is not shorter because you knocked louder.
\end{enumerate}

Composing the two layers (proactive bucket, reactive retry policy) yields a
client that stays under the limit in steady state and recovers gracefully when
it mis-estimates.

\subsection{Pagination and Streaming}
\label{sec:ch02-pagination}

APIs paginate to bound response size; consumers must abstract the mechanics
behind a lazy iterator so call sites write \texttt{for part in
client.parts():} and nothing more. The three dominant patterns
(Chapter~8 designs them; here we consume them):

\begin{description}
  \item[Offset/limit] --- \texttt{?offset=100\&limit=50}. Simple, stateless,
        but unstable under concurrent writes (rows shift between pages) and
        increasingly expensive for deep offsets. Termination: a short page or
        an empty one.
  \item[Cursor] --- \texttt{?cursor=eyJpZCI6...}. Opaque, stable, cheap at
        depth. The cursor must be treated as an opaque token --- never parsed,
        never fabricated. Termination: absent \texttt{next\_cursor}.
  \item[RFC 8288 Link headers] ---
        \texttt{Link: <...?page=3>; rel="next"}. Self-describing and
        cache-friendly; termination is the absence of \texttt{rel="next"}.
\end{description}

For large single payloads (catalog exports, firmware images), do not buffer:
\texttt{client.stream(...)} plus \texttt{response.iter\_bytes()} yields chunks
as they arrive, keeping memory $O(\text{chunk size})$ instead of
$O(\text{payload size})$. Three production details: (1) \texttt{httpx}
transparently negotiates and applies \texttt{gzip}/\texttt{br} decompression
when the codecs are installed --- verify \texttt{Content-Encoding} handling
rather than assuming it; (2) progress reporting must be driven by bytes
received against \texttt{Content-Length}, tolerating its absence
(chunked responses have none); (3) cancellation must propagate --- closing
the stream or cancelling the task must release the connection promptly.

\subsection{An Error Taxonomy for Consumers}
\label{sec:ch02-errors}

``The API call failed'' is not an error message; it is a shrug. Actionable
consumers classify failure along three orthogonal layers, because each layer
has a different owner and a different remedy:

\begin{table}[h]
  \centering
  \small
  \begin{tabularx}{\textwidth}{l X X X}
    \toprule
    \textbf{Layer} & \textbf{Examples} & \textbf{Meaning} & \textbf{Action} \\
    \midrule
    Transport & ConnectTimeout, ReadTimeout, ConnectError, pool exhaustion &
    Request never completed, or response never arrived & Retry (if
    idempotent); check network, DNS, TLS; alert on sustained rate \\
    Status (HTTP) & 4xx, 5xx after a complete exchange &
    Server answered and refused/failed & 4xx: fix the request (deterministic).
    5xx/429: retry per policy; page if sustained \\
    Contract / parse & JSON decode failure, schema violation, missing field &
    200 OK but the payload broke the contract & Never retry. Raise loudly;
    version skew or a provider defect --- treat as a bug \\
    \bottomrule
  \end{tabularx}
  \caption{Consumer error taxonomy. Each layer maps to a distinct exception
  type and remediation path.}
  \label{tab:ch02-error-taxonomy}
\end{table}

We therefore define a small exception hierarchy --- \texttt{ApiClientError}
at the root with \texttt{ApiTransportError}, \texttt{ApiStatusError}
(subclassed as \texttt{ApiClientError4xx}, \texttt{ApiServerError5xx},
\texttt{ApiRateLimited} carrying \texttt{retry\_after}), and
\texttt{ApiContractError} --- so callers can \texttt{except} precisely the
category they can act on. Kleppmann's systems lens
\cite{kleppmann2017ddia} generalizes the point: at a network boundary,
distinguish ``the request failed,'' ``the request succeeded,'' and
\emph{``we cannot know''} --- and design each path explicitly.

%% ---------------------------------------------------------------------------
%% 4. Production-Ready Code Implementation
%% ---------------------------------------------------------------------------

\section{Production-Ready Code Implementation}
\label{sec:ch02-code}

All listings target Python 3.11+, are fully type-annotated, and run offline.
The demonstration server is the local stub in
Listing~\ref{lst:ch02-stub}; every test uses \texttt{httpx.MockTransport}.
Run them with \texttt{python} (listings) and \texttt{pytest -q} (test suite)
from a virtual environment configured per
Section~\ref{sec:ch02-dependencies}.

\subsection{Baseline: A stdlib urllib Consumer}
\label{sec:ch02-urllib}

Listing~\ref{lst:ch02-urllib} is the deliberate ``before'' picture: correct,
dependency-free, and studded with the sharp edges catalogued in
Section~\ref{sec:ch02-evolution}. Read it and enumerate what production would
demand that it lacks.

\begin{examplebox}
\textbf{Reading the baseline.} As you read Listing~\ref{lst:ch02-urllib},
annotate: Where are the timeout dimensions? Where is keep-alive? What happens
on a 429? Where would authentication live? What is logged? The gaps you find
are the design brief for the httpx client that follows.
\end{examplebox}

\begin{minted}[caption={\texttt{ch02\_urllib\_baseline.py} --- stdlib baseline.},label={lst:ch02-urllib}]{python}
"""ch02_urllib_baseline.py -- stdlib urllib consumer (baseline, with sharp edges).

Run against the local stub:  python ch02_stub_server.py   (in another shell)
Then:                        python ch02_urllib_baseline.py
"""
from __future__ import annotations

import json
import logging
import urllib.error
import urllib.request
from dataclasses import dataclass
from typing import Any

logger = logging.getLogger("ch02.urllib_baseline")

BASE_URL = "http://127.0.0.1:8765"  # local stub only -- never the internet


@dataclass(frozen=True)
class Part:
    part_id: str
    name: str
    unit_price_cents: int


def fetch_part(part_id: str, timeout_s: float = 5.0) -> Part:
    """Fetch one part with urllib. NOTE the single, flat timeout.

    Sharp edges on display:
      * one timeout covers DNS+connect+read; no write/pool dimension;
      * no session -> no keep-alive: DNS/TCP/TLS are re-paid per call;
      * HTTPError is raised for 4xx/5xx AND doubles as a response object;
      * JSON, status classification, and retry are all manual.
    """
    url = f"{BASE_URL}/parts/{part_id}"
    req = urllib.request.Request(url, headers={"Accept": "application/json"})
    try:
        with urllib.request.urlopen(req, timeout=timeout_s) as resp:
            status: int = resp.status
            raw: bytes = resp.read()  # unbounded: a hostile peer can exhaust RAM
    except urllib.error.HTTPError as exc:
        # HTTPError is simultaneously the exception and the response.
        logger.error("HTTP %s for %s", exc.code, url)
        raise RuntimeError(f"status {exc.code}") from exc
    except urllib.error.URLError as exc:
        # Transport failure (DNS, refused connection, timeout -> socket.timeout
        # wrapped inside .reason). Nothing distinguishes the sub-causes here.
        logger.error("transport failure for %s: %s", url, exc.reason)
        raise RuntimeError(f"transport: {exc.reason}") from exc

    if status != 200:
        raise RuntimeError(f"unexpected status {status}")
    payload: dict[str, Any] = json.loads(raw.decode("utf-8"))
    return Part(
        part_id=str(payload["part_id"]),
        name=str(payload["name"]),
        unit_price_cents=int(payload["unit_price_cents"]),
    )


def main() -> None:
    logging.basicConfig(level=logging.INFO, format="%(levelname)s %(name)s: %(message)s")
    part = fetch_part("NR-1001")
    logger.info("got part: %s", part)


if __name__ == "__main__":
    main()
\end{minted}

\subsection{A Production-Grade httpx Client}
\label{sec:ch02-httpxclient}

Listing~\ref{lst:ch02-httpx} implements the pipeline of
Figure~\ref{fig:ch02-lifecycle}: four-dimension timeouts, bounded pooling,
single-point auth injection, event-hook telemetry, the error taxonomy of
Table~\ref{tab:ch02-error-taxonomy}, and transport injection so the identical
object runs against production URLs or an in-memory mock.

\begin{minted}[caption={\texttt{ch02\_httpx\_client.py} --- production client core.},label={lst:ch02-httpx}]{python}
"""ch02_httpx_client.py -- production-grade httpx client for the
Northwind Robotics Parts API.

Design contract:
  * explicit 4-dimension timeouts (connect/read/write/pool); nothing infinite;
  * bounded connection pool with keep-alive; one long-lived client instance;
  * auth injected in exactly one place, never logged;
  * typed error taxonomy: transport vs status vs contract;
  * telemetry via event hooks (latency + status, header-safe logging);
  * transport injection (httpx.MockTransport) for hermetic offline tests.
"""
from __future__ import annotations

import logging
import time
from typing import Any

import httpx

logger = logging.getLogger("ch02.httpx_client")

RETRYABLE_STATUSES: frozenset[int] = frozenset({408, 425, 429, 500, 502, 503, 504})


# ---- Error taxonomy ---------------------------------------------------------

class ApiClientError(Exception):
    """Root of all errors raised by this client."""


class ApiTransportError(ApiClientError):
    """The request never completed: DNS, TCP, TLS, timeout, pool exhaustion."""


class ApiStatusError(ApiClientError):
    """The server answered with a non-2xx status."""

    def __init__(self, message: str, *, status_code: int) -> None:
        super().__init__(message)
        self.status_code = status_code


class ApiClientError4xx(ApiStatusError):
    """Deterministic caller fault; retrying as-is is futile."""


class ApiServerError5xx(ApiStatusError):
    """Transient server fault; a candidate for retry per policy."""


class ApiRateLimited(ApiStatusError):
    """429: carries the server-advertised delay when present."""

    def __init__(self, message: str, *, status_code: int, retry_after: float | None) -> None:
        super().__init__(message, status_code=status_code)
        self.retry_after = retry_after


class ApiContractError(ApiClientError):
    """2xx received, but the payload violates the expected contract."""


# ---- Telemetry hooks ---------------------------------------------------------

def _log_request(request: httpx.Request) -> None:
    # NEVER log request.headers: Authorization lives there.
    logger.info("request %s %s", request.method, request.url.path)


def _log_response(response: httpx.Response) -> None:
    # Latency is timed in _request() with time.monotonic(): response.elapsed
    # is only populated after the body is read, which a hook must not force
    # (it would buffer streaming responses).
    logger.info(
        "response %s %s -> %d",
        response.request.method, response.request.url.path, response.status_code,
    )


# ---- The client ---------------------------------------------------------------

class NorthwindPartsClient:
    """Typed, resilient consumer for the Northwind Robotics Parts API."""

    def __init__(
        self,
        base_url: str,
        api_key: str,
        *,
        transport: httpx.BaseTransport | None = None,  # inject MockTransport in tests
        timeout: httpx.Timeout | None = None,
        limits: httpx.Limits | None = None,
    ) -> None:
        if not base_url:
            raise ValueError("base_url must be non-empty")
        if not api_key:
            raise ValueError("api_key must be non-empty")
        self._client = httpx.Client(
            base_url=base_url.rstrip("/"),
            # Auth lives HERE and nowhere else.
            headers={
                "Authorization": f"Bearer {api_key}",
                "Accept": "application/json",
                "User-Agent": "northwind-parts-sdk/1.0 (ops@northwind-robotics.example)",
            },
            timeout=timeout or httpx.Timeout(
                connect=2.0, read=10.0, write=10.0, pool=2.0
            ),
            limits=limits or httpx.Limits(
                max_connections=20, max_keepalive_connections=10,
                keepalive_expiry=30.0,
            ),
            event_hooks={"request": [_log_request], "response": [_log_response]},
            transport=transport,
        )

    # Context-manager protocol: deterministic socket lifecycle.
    def __enter__(self) -> "NorthwindPartsClient":
        return self

    def __exit__(self, *exc_info: object) -> None:
        self.close()

    def close(self) -> None:
        self._client.close()

    # Core request path ------------------------------------------------------

    def _request(self, method: str, path: str, **kwargs: Any) -> httpx.Response:
        started = time.monotonic()
        try:
            response = self._client.request(method, path, **kwargs)
        except httpx.TimeoutException as exc:
            kind = type(exc).__name__  # ConnectTimeout / ReadTimeout / PoolTimeout / ...
            logger.warning("timeout (%s) on %s %s", kind, method, path)
            raise ApiTransportError(f"{kind} on {method} {path}") from exc
        except httpx.TransportError as exc:
            logger.warning("transport error on %s %s: %s", method, path, exc)
            raise ApiTransportError(f"transport error on {method} {path}: {exc}") from exc
        finally:
            logger.debug(
                "call %s %s took %.1fms",
                method, path, (time.monotonic() - started) * 1000.0,
            )
        self._raise_for_status(response)
        return response

    @staticmethod
    def _raise_for_status(response: httpx.Response) -> None:
        code = response.status_code
        if code < 400:
            return
        message = f"{response.request.method} {response.request.url.path} -> {code}"
        if code == 429:
            retry_after = _parse_retry_after(response.headers.get("Retry-After"))
            raise ApiRateLimited(message, status_code=code, retry_after=retry_after)
        if code >= 500:
            raise ApiServerError5xx(message, status_code=code)
        raise ApiClientError4xx(message, status_code=code)

    # Public API --------------------------------------------------------------

    def get_json(self, path: str, **kwargs: Any) -> dict[str, Any]:
        response = self._request("GET", path, **kwargs)
        try:
            payload = response.json()
        except ValueError as exc:  # 200 OK but not JSON: contract breach.
            raise ApiContractError(f"{path}: invalid JSON body") from exc
        if not isinstance(payload, dict):
            raise ApiContractError(f"{path}: expected JSON object, got {type(payload).__name__}")
        return payload

    def get_part(self, part_id: str) -> dict[str, Any]:
        if not part_id:
            raise ValueError("part_id must be non-empty")
        return self.get_json(f"/parts/{part_id}")


def _parse_retry_after(value: str | None, *, ceiling_s: float = 60.0) -> float | None:
    """Parse a delta-seconds Retry-After header, clamped to a ceiling."""
    if value is None:
        return None
    try:
        seconds = float(value)  # delta-seconds form; HTTP-date form omitted for brevity
    except ValueError:
        return None
    return max(0.0, min(seconds, ceiling_s))


def main() -> None:
    logging.basicConfig(level=logging.INFO, format="%(levelname)s %(name)s: %(message)s")
    # The demo key is synthetic; the stub accepts any non-empty bearer token.
    with NorthwindPartsClient("http://127.0.0.1:8765", api_key="demo-key-nr-001") as client:
        part = client.get_part("NR-1001")
        logger.info("part: %s", part)


if __name__ == "__main__":
    main()
\end{minted}

\subsection{Retry Engineering with tenacity}
\label{sec:ch02-retry-code}

Listing~\ref{lst:ch02-retry} layers retry policy onto the client using
\texttt{tenacity}. The policy encodes exactly the rules of
Section~\ref{sec:ch02-retries}: full jitter via
\texttt{wait\_random\_exponential} (mean $b/2$ doubling to the cap, matching
Equation~\ref{eq:ch02-jitter-total}), \texttt{Retry-After} honored and clamped
when a 429 advertises it, a combined attempt-plus-elapsed budget, and an
idempotency gate that refuses to replay non-idempotent methods lacking an
idempotency key. Determinism for tests is achieved by injecting the random
source and a no-op sleep.

\begin{minted}[caption={\texttt{ch02\_retry.py} --- jitter, Retry-After, idempotency gate.},label={lst:ch02-retry}]{python}
"""ch02_retry.py -- retry wrapper: full jitter + Retry-After + idempotency gate.

Policy (matches Section 3.5):
  * retry on ApiTransportError, ApiServerError5xx, ApiRateLimited;
  * full jitter: wait ~ Uniform(0, min(cap, base * 2**(n-1))), base=0.5s, cap=8s;
  * on 429 with Retry-After, the server-advertised delay OVERRIDES the jitter;
  * budget: at most 4 attempts AND at most 20s total elapsed;
  * idempotency gate: non-idempotent methods retry only with an
    Idempotency-Key header;
  * deterministic tests: inject `rng` and `sleep`.
"""
from __future__ import annotations

import logging
import random
from collections.abc import Callable

import httpx
from tenacity import (
    RetryCallState,
    Retrying,
    retry_if_exception_type,
    stop_after_attempt,
    stop_after_delay,
)

from ch02_httpx_client import (
    ApiRateLimited,
    ApiServerError5xx,
    ApiTransportError,
    NorthwindPartsClient,
)

logger = logging.getLogger("ch02.retry")

BASE_DELAY_S = 0.5
CAP_DELAY_S = 8.0
MAX_ATTEMPTS = 4
MAX_ELAPSED_S = 20.0

SAFE_METHODS: frozenset[str] = frozenset({"GET", "HEAD", "OPTIONS", "PUT", "DELETE"})
RETRYABLE = (ApiTransportError, ApiServerError5xx, ApiRateLimited)


def wait_policy(state: RetryCallState, *, rng: random.Random) -> float:
    """Retry-After from a 429 wins; otherwise full jitter on exponential backoff.

    tenacity calls this between attempts with the just-failed attempt's state,
    so the server-advertised delay is read straight off the raised exception.
    """
    exc = state.outcome.exception() if state.outcome is not None else None
    if isinstance(exc, ApiRateLimited) and exc.retry_after is not None:
        logger.info("429 received: honoring Retry-After of %.1fs", exc.retry_after)
        return exc.retry_after
    upper = min(CAP_DELAY_S, BASE_DELAY_S * (2 ** (state.attempt_number - 1)))
    delay = rng.uniform(0.0, upper)
    logger.info("attempt %d failed; full-jitter backoff %.2fs",
                state.attempt_number, delay)
    return delay


class ResilientPartsClient:
    """NorthwindPartsClient + retry policy. Retryable failures only."""

    def __init__(
        self,
        inner: NorthwindPartsClient,
        *,
        rng: random.Random | None = None,
        sleep: Callable[[float], None] | None = None,
    ) -> None:
        self._inner = inner
        # Seeded RNG and injectable sleep keep tests deterministic & instant.
        self._rng = rng or random.Random(1337)
        self._sleep = sleep  # None -> tenacity's real sleep (production)

    def request_with_retry(
        self,
        method: str,
        path: str,
        *,
        idempotency_key: str | None = None,
        **kwargs: object,
    ) -> httpx.Response:
        method = method.upper()
        if method not in SAFE_METHODS and not idempotency_key:
            raise ValueError(
                f"refusing to auto-retry non-idempotent {method} without "
                "an idempotency key"
            )
        if idempotency_key:
            headers = dict(kwargs.pop("headers", {}) or {})
            headers["Idempotency-Key"] = idempotency_key
            kwargs["headers"] = headers

        retrying = Retrying(
            retry=retry_if_exception_type(RETRYABLE),
            wait=lambda state: wait_policy(state, rng=self._rng),
            stop=stop_after_attempt(MAX_ATTEMPTS) | stop_after_delay(MAX_ELAPSED_S),
            sleep=self._sleep if self._sleep is not None else _real_sleep,
            reraise=True,
        )
        return retrying(self._inner._request, method, path, **kwargs)


def _real_sleep(seconds: float) -> None:
    import time

    time.sleep(seconds)


def main() -> None:
    logging.basicConfig(level=logging.INFO, format="%(levelname)s %(name)s: %(message)s")
    inner = NorthwindPartsClient("http://127.0.0.1:8765", api_key="demo-key-nr-001")
    resilient = ResilientPartsClient(inner)
    try:
        response = resilient.request_with_retry("GET", "/parts/NR-1001")
        logger.info("status=%d body=%s", response.status_code, response.text[:80])
    finally:
        inner.close()


if __name__ == "__main__":
    main()
\end{minted}

\subsection{Pagination Iterators and a Streaming Downloader}
\label{sec:ch02-pagination-code}

Listing~\ref{lst:ch02-pagination} hides pagination mechanics behind lazy
typed iterators (cursor and link-header variants) and adds a streaming
downloader with byte-level progress and prompt cancellation.

\begin{minted}[caption={\texttt{ch02\_pagination\_stream.py} --- iterators and streaming.},label={lst:ch02-pagination}]{python}
"""ch02_pagination_stream.py -- typed pagination iterators + streaming download.

Cursor iterator: follows opaque `next_cursor` tokens; never parses them.
Link iterator: follows RFC 8288 `Link: <url>; rel="next"` headers.
Downloader:    streams via iter_bytes (bounded memory), reports progress,
               and honours cancellation by closing the response promptly.
"""
from __future__ import annotations

import logging
from collections.abc import Callable, Iterator
from pathlib import Path
from typing import Any

import httpx

from ch02_httpx_client import ApiContractError, NorthwindPartsClient

logger = logging.getLogger("ch02.pagination")

CHUNK_SIZE = 64 * 1024  # 64 KiB: memory stays O(chunk), never O(payload)


def iter_parts_cursor(
    client: NorthwindPartsClient,
    *,
    page_size: int = 50,
    max_pages: int = 1000,  # runaway-guard: a broken server must not loop us forever
) -> Iterator[dict[str, Any]]:
    """Yield every part from a cursor-paginated /parts endpoint."""
    cursor: str | None = None
    for _ in range(max_pages):
        params: dict[str, Any] = {"limit": page_size}
        if cursor is not None:
            params["cursor"] = cursor  # opaque token: pass through, never parse
        page = client.get_json("/parts", params=params)
        items = page.get("items")
        if not isinstance(items, list):
            raise ApiContractError("/parts: 'items' must be a list")
        yield from items
        cursor = page.get("next_cursor")
        if not cursor:  # termination: absent next_cursor
            return
    raise ApiContractError(f"/parts: exceeded {max_pages} pages (server bug?)")


def iter_parts_link_header(
    client: NorthwindPartsClient, *, max_pages: int = 1000
) -> Iterator[dict[str, Any]]:
    """Yield every part by following RFC 8288 Link: rel="next" headers."""
    response = client._request("GET", "/parts-link")
    for _ in range(max_pages):
        try:
            page: dict[str, Any] = response.json()
        except ValueError as exc:
            raise ApiContractError("/parts-link: invalid JSON body") from exc
        items = page.get("items")
        if not isinstance(items, list):
            raise ApiContractError("/parts-link: 'items' must be a list")
        yield from items
        next_url = _parse_link_next(response.headers.get("Link"))
        if next_url is None:  # termination: no rel="next"
            return
        response = client._request("GET", next_url)
    raise ApiContractError(f"/parts-link: exceeded {max_pages} pages")


def _parse_link_next(header: str | None) -> str | None:
    """Extract the rel="next" target from an RFC 8288 Link header."""
    if not header:
        return None
    for segment in header.split(","):
        parts = segment.split(";")
        if len(parts) == 2 and parts[1].strip() == 'rel="next"':
            return parts[0].strip().strip("<>")
    return None


def stream_download(
    client: NorthwindPartsClient,
    path: str,
    destination: Path,
    *,
    on_progress: Callable[[int, int | None], None] | None = None,
    cancel: Callable[[], bool] = lambda: False,
) -> int:
    """Stream `path` to `destination`; return bytes written.

    Memory stays O(CHUNK_SIZE). Cancellation is cooperative: when cancel()
    returns True we abort, close the stream, and delete the partial file so
    no corrupt artifact survives.
    """
    received = 0
    tmp = destination.with_suffix(destination.suffix + ".part")
    try:
        with client._client.stream("GET", path) as response:
            client._raise_for_status(response)
            total_header = response.headers.get("Content-Length")
            total = int(total_header) if total_header else None
            with tmp.open("wb") as fh:
                for chunk in response.iter_bytes(CHUNK_SIZE):
                    if cancel():
                        raise InterruptedError("download cancelled by caller")
                    fh.write(chunk)
                    received += len(chunk)
                    if on_progress:
                        on_progress(received, total)
        tmp.replace(destination)  # atomic rename: all-or-nothing artifact
        return received
    except BaseException:
        tmp.unlink(missing_ok=True)  # never leave a partial file behind
        raise


def main() -> None:
    logging.basicConfig(level=logging.INFO, format="%(levelname)s %(name)s: %(message)s")
    with NorthwindPartsClient("http://127.0.0.1:8765", api_key="demo-key-nr-001") as client:
        count = sum(1 for _ in iter_parts_cursor(client, page_size=2))
        logger.info("cursor-paginated parts: %d", count)
        written = stream_download(
            client, "/catalog/export",
            Path("catalog_export.jsonl"),
            on_progress=lambda done, total: logger.info("progress: %d/%s bytes", done, total),
        )
        logger.info("downloaded %d bytes", written)


if __name__ == "__main__":
    main()
\end{minted}

\subsection{Hermetic Offline Tests with httpx.MockTransport}
\label{sec:ch02-tests}

The final listing is the discipline that makes everything above trustworthy:
a fault-matrix suite that exercises the full client stack --- timeouts,
pooling, hooks, retry policy, taxonomy --- with zero sockets, zero wall-clock
dependence, and seeded randomness. \texttt{MockTransport} routes requests to
an in-process handler while preserving the entire \texttt{httpx} pipeline
\cite{httpxdocs}; the injected \texttt{sleep} records delays instead of
consuming real time.

\begin{minted}[caption={\texttt{test\_ch02\_client.py} --- offline fault-matrix test suite.},label={lst:ch02-tests}]{python}
"""test_ch02_client.py -- hermetic tests: no network, no wall clock, seeded RNG.

Run:  pytest -q
Simulates: 500s then success, 429 with Retry-After, deterministic 4xx,
transport timeouts, slow responses (mocked), contract breaches, and the
idempotency gate.
"""
from __future__ import annotations

import random

import httpx
import pytest

from ch02_httpx_client import (
    ApiClientError4xx,
    ApiContractError,
    ApiRateLimited,
    ApiServerError5xx,
    ApiTransportError,
    NorthwindPartsClient,
)
from ch02_retry import ResilientPartsClient

BASE_URL = "http://stub.local"  # never resolves: MockTransport intercepts


def make_client(handler, **kwargs) -> NorthwindPartsClient:
    return NorthwindPartsClient(
        BASE_URL, api_key="test-key",
        transport=httpx.MockTransport(handler), **kwargs,
    )


# --- Status handling & taxonomy ---------------------------------------------

def test_get_part_success() -> None:
    def handler(request: httpx.Request) -> httpx.Response:
        assert request.headers["Authorization"] == "Bearer test-key"
        return httpx.Response(200, json={"part_id": "NR-1001", "name": "Servo"})
    with make_client(handler) as client:
        assert client.get_part("NR-1001")["name"] == "Servo"


def test_404_maps_to_typed_client_error() -> None:
    with make_client(lambda r: httpx.Response(404, json={"error": "no such part"})) as c:
        with pytest.raises(ApiClientError4xx) as exc_info:
            c.get_part("NOPE-1")
        assert exc_info.value.status_code == 404


def test_500_maps_to_server_error() -> None:
    with make_client(lambda r: httpx.Response(503)) as c:
        with pytest.raises(ApiServerError5xx):
            c.get_part("NR-1001")


def test_429_carries_retry_after() -> None:
    def handler(request: httpx.Request) -> httpx.Response:
        return httpx.Response(429, headers={"Retry-After": "7"})
    with make_client(handler) as c:
        with pytest.raises(ApiRateLimited) as exc_info:
            c.get_part("NR-1001")
        assert exc_info.value.retry_after == 7.0


def test_transport_timeout_maps_to_transport_error() -> None:
    def handler(request: httpx.Request) -> httpx.Response:
        raise httpx.ReadTimeout("simulated hung peer", request=request)
    with make_client(handler) as c:
        with pytest.raises(ApiTransportError):
            c.get_part("NR-1001")


def test_200_with_broken_body_is_contract_error() -> None:
    with make_client(lambda r: httpx.Response(200, content=b"<html>oops</html>")) as c:
        with pytest.raises(ApiContractError):
            c.get_part("NR-1001")


# --- Retry policy --------------------------------------------------------------

def test_retry_recovers_from_transient_500s() -> None:
    calls = {"n": 0}
    def handler(request: httpx.Request) -> httpx.Response:
        calls["n"] += 1
        if calls["n"] < 3:
            return httpx.Response(500)
        return httpx.Response(200, json={"part_id": "NR-1001"})
    inner = make_client(handler)
    sleeps: list[float] = []
    resilient = ResilientPartsClient(
        inner, rng=random.Random(42), sleep=sleeps.append
    )
    response = resilient.request_with_retry("GET", "/parts/NR-1001")
    assert response.status_code == 200
    assert calls["n"] == 3
    assert len(sleeps) == 2  # two retries, two jittered sleeps, zero wall time


def test_retry_honors_retry_after_over_jitter() -> None:
    calls = {"n": 0}
    def handler(request: httpx.Request) -> httpx.Response:
        calls["n"] += 1
        if calls["n"] == 1:
            return httpx.Response(429, headers={"Retry-After": "5"})
        return httpx.Response(200, json={"ok": True})
    inner = make_client(handler)
    sleeps: list[float] = []
    resilient = ResilientPartsClient(
        inner, rng=random.Random(42), sleep=sleeps.append
    )
    resilient.request_with_retry("GET", "/parts/NR-1001")
    assert sleeps == [5.0]  # server-advertised delay, not jitter


def test_4xx_is_never_retried() -> None:
    calls = {"n": 0}
    def handler(request: httpx.Request) -> httpx.Response:
        calls["n"] += 1
        return httpx.Response(422, json={"detail": "invalid"})
    inner = make_client(handler)
    resilient = ResilientPartsClient(inner, rng=random.Random(42),
                                     sleep=lambda s: None)
    with pytest.raises(ApiClientError4xx):
        resilient.request_with_retry("GET", "/parts/NR-1001")
    assert calls["n"] == 1


def test_post_without_idempotency_key_is_refused() -> None:
    inner = make_client(lambda r: httpx.Response(200, json={}))
    resilient = ResilientPartsClient(inner)
    with pytest.raises(ValueError, match="idempotency"):
        resilient.request_with_retry("POST", "/orders", json={"qty": 1})


def test_post_with_idempotency_key_retries_and_sends_header() -> None:
    seen: list[str | None] = []
    calls = {"n": 0}
    def handler(request: httpx.Request) -> httpx.Response:
        seen.append(request.headers.get("Idempotency-Key"))
        calls["n"] += 1
        if calls["n"] == 1:
            return httpx.Response(503)
        return httpx.Response(201, json={"order_id": "ORD-9"})
    inner = make_client(handler)
    resilient = ResilientPartsClient(inner, rng=random.Random(42),
                                     sleep=lambda s: None)
    response = resilient.request_with_retry(
        "POST", "/orders", idempotency_key="key-abc-123", json={"qty": 1}
    )
    assert response.status_code == 201
    assert seen == ["key-abc-123", "key-abc-123"]  # key stable across replays
\end{minted}

Note what the suite does \emph{not} do: open a socket, sleep in real time,
assert on wall-clock, or depend on test ordering. Those four absences are the
definition of a deterministic, offline-first consumer test suite.

%% ---------------------------------------------------------------------------
%% 5. Common Pitfalls & Anti-Patterns
%% ---------------------------------------------------------------------------

\section{Common Pitfalls \& Anti-Patterns}
\label{sec:ch02-pitfalls}

Each anti-pattern below is a real production failure mode, stated with its
mechanism and its fix.

\subsection{Missing Timeouts}

\textbf{Anti-pattern:} \texttt{requests.get(url)} --- infinite wait by
default. \textbf{Mechanism:} a hung peer holds a worker forever; under load
the pool exhausts and the healthy caller falls over (the cascading failure of
Section~\ref{sec:ch02-warstory}). \textbf{Fix:} explicit four-dimension
timeouts on the client, tighter or looser per call only with justification.
Treat any timeout-free call found in review as a defect, full stop.

\subsection{Retry Storms}

\textbf{Anti-pattern:} unbounded \texttt{while True: try again immediately},
or fixed-interval retries with no cap, no jitter, no budget.
\textbf{Mechanism:} when the dependency degrades, every client retries in
lockstep; offered load multiplies exactly when capacity shrinks; the outage
deepens and self-sustains. \textbf{Fix:} capped attempts \emph{plus} elapsed
budget, full jitter, \texttt{Retry-After} honored, and a retry budget
(Section~\ref{sec:ch02-retries}); circuit breaking for sustained failure
(Chapter~15).

\subsection{Swallowing Exceptions}

\textbf{Anti-pattern:} \texttt{except Exception: pass}, or logging and
returning \texttt{None} from a function whose contract promises a resource.
\textbf{Mechanism:} failures become silent data corruption --- a ``missing
part'' that was really a \texttt{503}; downstream code manufactures defaults;
the root cause is invisible to operators. \textbf{Fix:} catch specific types
(the taxonomy of Section~\ref{sec:ch02-errors}), add context, re-raise; let
only code that can \emph{act} on a failure catch it.

\subsection{Parsing Before the Status Check}

\textbf{Anti-pattern:} \texttt{resp.json()} unconditionally.
\textbf{Mechanism:} a \texttt{503} with an HTML body from a proxy surfaces as
\texttt{JSONDecodeError}, three frames from the truth; on-call engineers debug
serialization when the real story is capacity. \textbf{Fix:} classify the
status first (our \texttt{\_raise\_for\_status}), then parse; a 2xx with an
unparseable body is a distinct \texttt{ApiContractError}.

\subsection{Leaking Credentials into Logs}

\textbf{Anti-pattern:} logging \texttt{request.headers}, full URLs with query
tokens, or response bodies on the error path. \textbf{Mechanism:} bearer
tokens and API keys land in log aggregators with broad read access and long
retention --- a credential-leak incident from a debugging convenience.
\textbf{Fix:} log method, path, status, latency; never headers, bodies, or
auth material (see \texttt{\_log\_request} in Listing~\ref{lst:ch02-httpx}).
The twelve-factor config discipline \cite{twelvefactor} applies: credentials
arrive via environment, never via source, and never leave via logs.

\subsection{Unbounded Decompression and Body Reads}

\textbf{Anti-pattern:} \texttt{response.read()} / \texttt{response.text} on
untrusted or unbounded endpoints; accepting \texttt{gzip}/\texttt{br} without
size discipline. \textbf{Mechanism:} a 2~MB gzip bomb expands to gigabytes in
memory; a streaming endpoint with no \texttt{Content-Length} reads until RAM
is gone. \textbf{Fix:} stream with \texttt{iter\_bytes}
(Listing~\ref{lst:ch02-pagination}) for anything large or untrusted; enforce a
maximum decompressed size; write partial downloads to a \texttt{.part} file
and rename atomically on success.

\begin{warningbox}
\textbf{The compound anti-pattern.} These failures compose: no timeout +
aggressive retry + swallowed exception is a client that hangs forever, retries
forever, and reports nothing. Review consumers as systems, not as lines.
\end{warningbox}

%% ---------------------------------------------------------------------------
%% 6. Hands-On Production Lab & Real-World Case Study
%% ---------------------------------------------------------------------------

\section{Hands-On Production Lab \& Real-World Case Study}
\label{sec:ch02-lab}

\subsection{Lab Setup: The Northwind Robotics Parts API Stub}

The lab consumes a fictional supplier API --- the \emph{Northwind Robotics
Parts API} --- served entirely by a local stub on \texttt{127.0.0.1:8765}.
The stub (Listing~\ref{lst:ch02-stub}) uses only the standard library and
implements five behaviors: a part lookup, cursor pagination over a small
synthetic catalog, a \texttt{429 + Retry-After} fault for the retry demo, a
flaky endpoint that fails twice before succeeding, and a streaming catalog
export.

\begin{minted}[caption={\texttt{ch02\_stub\_server.py} --- local stub for the Northwind Robotics Parts API.},label={lst:ch02-stub}]{python}
"""ch02_stub_server.py -- offline stub for the Northwind Robotics Parts API.

Stdlib only. Deterministic. Listens on 127.0.0.1:8765 (localhost only).

Endpoints:
  GET /parts/{id}        -> 200 JSON part, or 404
  GET /parts?limit&cursor-> cursor-paginated catalog (5 synthetic parts)
  GET /parts-link        -> RFC 8288 Link-header pagination (2 pages)
  GET /rate-limited      -> 429 with Retry-After: 2 (always)
  GET /flaky             -> 500, 500, then 200 (per-process state)
  GET /catalog/export    -> streamed JSONL catalog

Run:  python ch02_stub_server.py
"""
from __future__ import annotations

import json
import logging
import threading
from http.server import BaseHTTPRequestHandler, ThreadingHTTPServer
from urllib.parse import parse_qs, urlparse

logger = logging.getLogger("ch02.stub")

# Synthetic Northwind Robotics catalog: fictional parts, fictional prices.
CATALOG: list[dict[str, object]] = [
    {"part_id": "NR-1001", "name": "Servo Motor SG-90", "unit_price_cents": 450},
    {"part_id": "NR-1002", "name": "LiDAR Unit LR-2", "unit_price_cents": 8999},
    {"part_id": "NR-1003", "name": "Chassis Plate AL-7075", "unit_price_cents": 3200},
    {"part_id": "NR-1004", "name": "Wheel Encoder WE-512", "unit_price_cents": 1250},
    {"part_id": "NR-1005", "name": "Battery Pack BP-4S", "unit_price_cents": 5400},
]

_flaky_calls = {"n": 0}
_state_lock = threading.Lock()


class StubHandler(BaseHTTPRequestHandler):
    server_version = "NorthwindStub/1.0"

    def _send_json(self, status: int, payload: dict[str, object],
                   extra_headers: dict[str, str] | None = None) -> None:
        body = json.dumps(payload).encode("utf-8")
        self.send_response(status)
        self.send_header("Content-Type", "application/json")
        self.send_header("Content-Length", str(len(body)))
        for key, value in (extra_headers or {}).items():
            self.send_header(key, value)
        self.end_headers()
        self.wfile.write(body)

    def do_GET(self) -> None:  # noqa: N802 (stdlib naming)
        parsed = urlparse(self.path)
        route, query = parsed.path, parse_qs(parsed.query)

        if route.startswith("/parts/"):
            part_id = route.rsplit("/", 1)[-1]
            part = next((p for p in CATALOG if p["part_id"] == part_id), None)
            if part is None:
                self._send_json(404, {"error": f"no such part: {part_id}"})
            else:
                self._send_json(200, part)
            return

        if route == "/parts":
            limit = int(query.get("limit", ["50"])[0])
            cursor = query.get("cursor", [None])[0]
            start = int(cursor) if cursor is not None else 0  # stub: cursor = index
            items = CATALOG[start:start + limit]
            next_index = start + limit
            payload: dict[str, object] = {"items": items}
            if next_index < len(CATALOG):
                payload["next_cursor"] = str(next_index)
            self._send_json(200, payload)
            return

        if route == "/parts-link":
            if query.get("page") == ["2"]:
                # Final page: no Link header -> iterator terminates.
                self._send_json(200, {"items": CATALOG[3:]})
            else:
                self._send_json(
                    200, {"items": CATALOG[:3]},
                    extra_headers={"Link": '</parts-link?page=2>; rel="next"'},
                )
            return

        if route == "/rate-limited":
            self._send_json(429, {"error": "slow down"},
                            extra_headers={"Retry-After": "2"})
            return

        if route == "/flaky":
            with _state_lock:
                _flaky_calls["n"] += 1
                n = _flaky_calls["n"]
            if n < 3:
                self._send_json(500, {"error": "transient blip"})
            else:
                self._send_json(200, {"part_id": "NR-1001", "ok": True})
            return

        if route == "/catalog/export":
            lines = "".join(json.dumps(p) + "\n" for p in CATALOG).encode("utf-8")
            self.send_response(200)
            self.send_header("Content-Type", "application/x-ndjson")
            self.send_header("Content-Length", str(len(lines)))
            self.end_headers()
            self.wfile.write(lines)
            return

        self._send_json(404, {"error": f"unknown route: {route}"})

    def log_message(self, fmt: str, *args: object) -> None:
        logger.info("stub: " + fmt, *args)


def main() -> None:
    logging.basicConfig(level=logging.INFO, format="%(levelname)s %(name)s: %(message)s")
    server = ThreadingHTTPServer(("127.0.0.1", 8765), StubHandler)
    logger.info("Northwind stub listening on http://127.0.0.1:8765 (Ctrl+C to stop)")
    try:
        server.serve_forever()
    except KeyboardInterrupt:
        server.shutdown()


if __name__ == "__main__":
    main()
\end{minted}

\subsection{Lab Steps}

Work from a project folder, e.g.\ \texttt{C:\textbackslash Training\textbackslash ch02\_lab},
containing the five listing files. All commands are PowerShell.

\begin{enumerate}
  \item \textbf{Environment.} Create and activate the virtual environment,
        then install the pinned dependencies:
\begin{minted}{bash}
python -m venv .venv
.\.venv\Scripts\Activate.ps1
pip install -r requirements.txt
\end{minted}
  \item \textbf{Start the stub} in a first terminal:
        \texttt{python ch02\_stub\_server.py}. Confirm it answers:
        \texttt{curl.exe http://127.0.0.1:8765/parts/NR-1001}.
  \item \textbf{Baseline.} In a second terminal run
        \texttt{python ch02\_urllib\_baseline.py}. Note in the stub's log
        that every request is a fresh connection (\texttt{Connection: close}
        semantics) --- the churn of Section~\ref{sec:ch02-pooling} made
        visible.
  \item \textbf{Production client.} Run \texttt{python ch02\_httpx\_client.py}.
        Observe the telemetry hooks emitting per-request latency lines.
  \item \textbf{Retry demo.} Point the demo at the flaky endpoint (edit
        \texttt{main()} to call \texttt{/flaky}), run
        \texttt{python ch02\_retry.py}, and watch two jittered retries
        precede success. Then call \texttt{/rate-limited} and confirm the
        client waits the server-advertised 2~seconds.
  \item \textbf{Pagination and streaming.} Run
        \texttt{python ch02\_pagination\_stream.py}; verify the iterator
        yields all five parts across three pages and that
        \texttt{catalog\_export.jsonl} contains five NDJSON lines.
  \item \textbf{Offline fault matrix.} Stop the stub and run
        \texttt{pytest -q}. All tests must pass with the stub down --- proof
        the suite is hermetic.
  \item \textbf{Break it on purpose.} Set the client's read timeout to
        \texttt{0.05} against the stub's \texttt{/flaky} endpoint, observe
        \texttt{ApiTransportError}, then restore. Change the stub's 429 to
        omit \texttt{Retry-After} and confirm the client falls back to full
        jitter.
\end{enumerate}

\begin{keyconceptbox}
\textbf{Exit criteria.} The lab is complete when: (1) the retry demo shows
server-advertised waits overriding jitter; (2) the pagination iterator
terminates without a \texttt{max\_pages} tripwire firing; (3) the pytest suite
passes with the network stack provably unused.
\end{keyconceptbox}

\subsection{War Story: The Missing Read Timeout and the Retry Storm}
\label{sec:ch02-warstory}

A fulfillment service (names fictionalized, physics real) called a warehouse
inventory API on every order. The call was a one-liner:
\texttt{requests.get(url)}. No timeout --- \texttt{requests} defaults to
infinite. A hand-rolled retry loop sat above it: \texttt{for attempt in
range(20): try: ...; time.sleep(1)}.

At 14:07 on a Tuesday, a bad deploy left the inventory service's worker
threads deadlocked on a database lock. The load balancer kept passing health
checks (the \texttt{/healthz} handler was a separate thread), so traffic kept
flowing into a process that accepted connections and then \emph{said nothing,
forever}. TCP was healthy; HTTP was catatonic. Only a read timeout can see
this failure --- and there was none.

The cascade ran exactly as Nygard's playbook predicts
\cite{nygard2018releaseit}. By 14:12 every fulfillment worker held a socket
awaiting bytes that would never come. By 14:15 the retry loops --- one-second
fixed interval, no jitter, twenty attempts --- had multiplied offered load on
the inventory service by roughly 20$\times$, ensuring that even when the lock
was manually killed at 14:22, the service drowned instantly in synchronized
retries: a self-inflicted retry storm. The circuit breaker that would have
shed the load did not exist. Checkout error rates peaked at 94\%; full
recovery took another 40 minutes after the root cause was fixed, because the
storm outlived its trigger.

The remediation diff was embarrassingly small: a four-dimension timeout
(connect 2~s, read 5~s, write 5~s, pool 1~s), three jittered retries with an
elapsed-time budget, \texttt{Retry-After} honored, and a fleet-wide retry
budget feeding a circuit breaker (Chapter~15). The lesson this chapter
encodes: \textbf{the cheap resilience work is done in the client, before the
incident, or it is done in the incident review, after.}

%% ---------------------------------------------------------------------------
%% 7. Self-Assessment Exercises & Deep-Dive Questions
%% ---------------------------------------------------------------------------

\section{Self-Assessment Exercises \& Deep-Dive Questions}

\subsection{Exercises}

\begin{enumerate}
  \item \textbf{Timeout audit.} Grep a codebase you know for
        \texttt{requests.get(}, \texttt{requests.post(}, and
        \texttt{urlopen(} lacking a \texttt{timeout=} argument. For each hit,
        state which of the four timeout dimensions is unbounded and what
        failure it permits.
  \item \textbf{Jitter analysis.} With $b = 0.5\,\text{s}$ and $N = 4$
        retries, compute (a) the deterministic exponential total wait,
        (b) the expected total wait under full jitter, and (c) the worst-case
        total wait under full jitter with cap $c = 8\,\text{s}$.
  \item \textbf{Idempotency triage.} For each operation, state whether a
        blind retry after a read timeout is safe, and why:
        \texttt{GET /parts/NR-1001}; \texttt{POST /orders};
        \texttt{PUT /parts/NR-1001/status} with body
        \texttt{\{"status": "reserved"\}}; \texttt{POST /payments} with an
        \texttt{Idempotency-Key} header.
  \item \textbf{Pool sizing.} A service runs 32 worker threads; each order
        makes at most 2 concurrent calls to one upstream. Propose
        \texttt{max\_connections} and \texttt{max\_keepalive\_connections}
        and justify both.
  \item \textbf{Extend the taxonomy.} Add \texttt{ApiAuthError} (401/403) to
        Listing~\ref{lst:ch02-httpx}, wire it into
        \texttt{\_raise\_for\_status}, and add a MockTransport test proving
        it is raised and never retried.
  \item \textbf{Offset walker.} The iterators in
        Listing~\ref{lst:ch02-pagination} cover cursor and link-header
        pagination. Implement \texttt{iter\_parts\_offset(client, limit)}
        for an offset/limit endpoint, with correct termination and a
        \texttt{max\_pages} tripwire, and test it against MockTransport.
\end{enumerate}

\subsection{Deep-Dive Questions}

\begin{enumerate}
  \item Why does a pool timeout exist separately from the read timeout? What
        distinct failure does each bound, and what does \texttt{PoolTimeout}
        tell you about your configuration that \texttt{ReadTimeout} never
        could?
  \item Full jitter halves expected total wait relative to deterministic
        backoff (Equation~\ref{eq:ch02-jitter-total}). Is that a latency cost
        or a latency saving for the caller, and why is it a fleet-level win
        regardless?
  \item A server sends \texttt{Retry-After: 120} but your retry budget caps
        elapsed time at 20~s. Reconcile the two policies. Which wins, and
        what should the client do when they conflict irreconcilably?
  \item Explain why \texttt{requests} cannot support HTTP/2 without a rewrite
        of its connection core, and what architectural property of
        \texttt{httpx} made HTTP/2 an \emph{option} rather than a fork.
  \item Under what conditions is cursor pagination strictly more correct than
        offset pagination for a catalog that changes while you read it?
        Construct the interleaving that corrupts an offset walk.
  \item Your MockTransport suite passes and production fails on TLS. What
        layer did the mock legitimately exclude, and what is the minimum
        additional test that covers it without internet access?
\end{enumerate}

%% ---------------------------------------------------------------------------
%% 8. Interview Q&A Vault
%% ---------------------------------------------------------------------------

\section{Interview Q\&A Vault}
\label{sec:ch02-interview}

Thirty-two questions ordered junior $\rightarrow$ staff. Answers are
deliberately compressed to what a strong candidate says out loud.

\subsection{Junior}

\begin{enumerate}
  \item \textbf{Q: What does an HTTP client library actually do for you?}\\
        \textit{A:} Manages the socket lifecycle (DNS, TCP, TLS), formats the
        request on the wire, parses the response, and --- in session-based
        libraries --- pools connections for reuse. Without it you are writing
        a state machine over sockets by hand.

  \item \textbf{Q: \texttt{urllib} vs.\ \texttt{requests} --- why does anyone
        still use \texttt{urllib}?}\\
        \textit{A:} Zero dependencies: it is always present in the stdlib, so
        installers, bootstrap scripts, and frozen single-file tools use it.
        For anything sustained, its flat timeout, missing session/pool, and
        manual JSON make \texttt{requests} or \texttt{httpx} the right call.

  \item \textbf{Q: What is keep-alive and why does it matter?}\\
        \textit{A:} Reusing one TCP/TLS connection across requests instead of
        reopening per request. It eliminates repeated handshake latency
        ($\sim$2 RTTs) and crypto CPU on every call --- often the single
        biggest client-side latency win.

  \item \textbf{Q: What is the default timeout of \texttt{requests.get()}?}\\
        \textit{A:} None --- it waits forever. That is a trap: the library's
        ergonomic default is a production hazard, which is why timeouts must
        always be explicit.

  \item \textbf{Q: A call returns status 500. What do you do?}\\
        \textit{A:} Classify it as a transient server fault: retry per policy
        (bounded, jittered, idempotency-aware). Never parse the body as if it
        were a success payload.

  \item \textbf{Q: A call returns 404. Retry it?}\\
        \textit{A:} No. 4xx is deterministic --- the same request fails
        identically. Fix the request (the URL, the ID, the payload), not the
        plumbing.

  \item \textbf{Q: What is JSON deserialization doing at the client
        boundary?}\\
        \textit{A:} Converting the response body bytes into Python objects.
        It belongs \emph{after} the status check, and its failure on a 2xx is
        a contract error, not a transport error.

  \item \textbf{Q: Why should a client send a \texttt{User-Agent}?}\\
        \textit{A:} Identifiability. When your traffic misbehaves, the
        provider can reach you instead of blocking an entire IP range --- and
        many providers contractually require it.
\end{enumerate}

\subsection{Mid-Level}

\begin{enumerate}
  \setcounter{enumi}{8}
  \item \textbf{Q: Why is a missing timeout a severity-1 risk?}\\
        \textit{A:} Because it converts a dependency's partial failure into
        your total failure: every worker blocks on a socket that never
        answers, the pool exhausts, and your healthy service stops serving.
        Unbounded waits are unbounded resource holds --- the textbook
        cascading failure \cite{nygard2018releaseit}.

  \item \textbf{Q: Name the four httpx timeout dimensions and what each
        bounds.}\\
        \textit{A:} connect (DNS/TCP/TLS establishment), read (waiting for
        response bytes, including mid-stream gaps), write (sending the
        request body), pool (waiting for a free connection in a saturated
        pool). Each bounds a distinct hang a single flat timeout cannot
        express.

  \item \textbf{Q: Connect timeout vs.\ read timeout --- which failure does
        each catch?}\\
        \textit{A:} Connect catches a dead, firewalled, or TLS-broken host.
        Read catches a live connection with a hung application behind it ---
        the failure TCP keepalive cannot see.

  \item \textbf{Q: What does \texttt{PoolTimeout} tell you?}\\
        \textit{A:} Your concurrency exceeds your pool: requests are queueing
        for connections. The fix is sizing (\texttt{max\_connections}) or
        concurrency control --- not a longer pool timeout, which just makes
        the queue quieter.

  \item \textbf{Q: Explain exponential backoff with full jitter.}\\
        \textit{A:} On attempt $n$, wait a \emph{random} value drawn
        uniformly from $[0, \min(c, b \cdot 2^{n-1})]$. Expected wait is
        $b \cdot 2^{n-2}$: exponential growth for server recovery, randomized
        phase so correlated clients stop colliding.

  \item \textbf{Q: Why decorrelated jitter?}\\
        \textit{A:} It samples the next wait from $[b, 3 \times
        \text{previous}]$ rather than the attempt index, so a client's
        schedule adapts to its own history. Both it and full jitter break
        lockstep; decorrelated jitter converges faster after short blips. The
        essential property either way is randomness --- without it, fleets
        retry in synchronized waves.

  \item \textbf{Q: What is a retry storm and how do you prevent it?}\\
        \textit{A:} Retries multiplying offered load on an already-degraded
        dependency, deepening the outage. Prevention: capped attempts,
        elapsed budgets, jitter, honoring \texttt{Retry-After}, a fleet-wide
        retry budget, and a circuit breaker for sustained failure.

  \item \textbf{Q: When is it safe to retry a POST?}\\
        \textit{A:} Only when the API accepts an idempotency key
        (server-side dedup) or the operation is naturally idempotent per its
        contract. Otherwise a read timeout is ambiguous --- the server may
        have succeeded --- and replaying risks duplicate side effects.

  \item \textbf{Q: requests vs.\ httpx trade-offs?}\\
        \textit{A:} \texttt{requests}: older, sync-only, HTTP/1.1-only,
        two-dimension timeouts, huge install base. \texttt{httpx}: same API
        shape plus sync \emph{and} async clients, optional HTTP/2,
        four-dimension timeouts, first-class MockTransport, full type hints.
        New code: httpx; existing sync codebases: requests is defensible.

  \item \textbf{Q: What does \texttt{Retry-After} do and where does it
        appear?}\\
        \textit{A:} Server-advertised minimum delay before retrying, defined
        by RFC~9110 \cite{rfc9110} for \texttt{429} and \texttt{503} as
        delta-seconds or an HTTP-date. It outranks any client-computed
        backoff because the server knows its own recovery timeline.

  \item \textbf{Q: How do you consume a paginated endpoint cleanly?}\\
        \textit{A:} Hide the mechanics behind a lazy iterator: the caller
        writes one \texttt{for} loop; the walker follows cursors or
        \texttt{Link: rel="next"} headers, treats cursors as opaque, and
        terminates on a missing next pointer --- with a page-count tripwire
        against server bugs.

  \item \textbf{Q: Why stream downloads instead of \texttt{response.read()}?}\\
        \textit{A:} Memory stays $O(\text{chunk})$ rather than
        $O(\text{payload})$; you get progress reporting, cooperative
        cancellation, and immunity to decompression bombs. Buffering is only
        acceptable for small, trusted, size-known payloads.
\end{enumerate}

\subsection{Senior}

\begin{enumerate}
  \setcounter{enumi}{20}
  \item \textbf{Q: Design the exception hierarchy for an SDK wrapper.}\\
        \textit{A:} One root (\texttt{ApiClientError}); three branches
        matching failure layers: transport (timeout/DNS/TLS --- retryable),
        status (split 4xx deterministic from 5xx transient; 429 carries
        \texttt{retry\_after}), and contract (2xx with a broken body ---
        never retry, alert). Callers \texttt{except} the layer they can act
        on.

  \item \textbf{Q: How do you test retry logic without slow, flaky tests?}\\
        \textit{A:} Inject everything time-like: \texttt{httpx.MockTransport}
        for the wire, a seeded \texttt{random.Random} for jitter, and an
        injected \texttt{sleep} that records durations instead of sleeping.
        Assert on call counts and recorded delays --- zero wall-clock
        assertions.

  \item \textbf{Q: What is a retry budget and why aren't per-request caps
        enough?}\\
        \textit{A:} Per-request caps bound one call; a budget bounds the
        fleet --- e.g., retries limited to 10\% of issued requests per
        minute. When a dependency fails broadly, every request hitting its
        per-request cap still multiplies total load; the budget converts that
        into controlled fail-fast.

  \item \textbf{Q: Your client sees \texttt{429} with \texttt{Retry-After:
        120} but an elapsed budget of 20~s. Design the resolution.}\\
        \textit{A:} The server delay wins within the budget; on conflict,
        fail the operation with \texttt{ApiRateLimited} carrying the delay so
        the \emph{caller} can requeue the work at the application layer.
        Never silently clamp to 20~s and retry into a still-limited server.

  \item \textbf{Q: How do event hooks differ from subclassing the client?}\\
        \textit{A:} Hooks are cross-cutting seams (logging, tracing, auth
        refresh) applied uniformly without touching the request path;
        subclassing embeds policy in the call path and fragments across
        overrides. Prefer composition (hooks, wrapped transports) over
        inheritance.

  \item \textbf{Q: A 2xx response contains HTML from a captive proxy. Which
        layer failed and how should the client react?}\\
        \textit{A:} Contract layer. \texttt{Content-Type} is wrong and JSON
        parsing fails; raise \texttt{ApiContractError}, alert, and never
        retry --- replaying cannot fix a middlebox interposing.

  \item \textbf{Q: Why is creating a new \texttt{httpx.Client} per request an
        anti-pattern?}\\
        \textit{A:} It discards the connection pool: every call re-pays DNS,
        TCP, and TLS, and sockets accumulate in \texttt{TIME\_WAIT} under
        load. One long-lived client per base URL, closed deterministically,
        is the correct shape.

  \item \textbf{Q: What telemetry belongs on the client hot path, and what
        must never appear in logs?}\\
        \textit{A:} Method, path, status class, latency, retry count,
        exception type. Never: \texttt{Authorization} headers, API keys,
        cookies, or full bodies --- credentials in log aggregation is a
        reportable leak.

  \item \textbf{Q: HTTP/2 changes what for a client?}\\
        \textit{A:} Multiplexing: many concurrent streams share one
        connection, eliminating HTTP/1.1 head-of-line blocking at the
        connection level and shrinking pool pressure. It does \emph{not}
        change timeout, retry, or idempotency semantics --- those are
        application-layer concerns.

  \item \textbf{Q: When would you choose async consumption, and what does it
        actually buy?}\\
        \textit{A:} When a single process must sustain high concurrent
        outbound calls (crawlers, fan-out aggregators). It buys massive
        connection concurrency without thread-per-connection memory and
        context-switch costs. It buys nothing for a sequential script ---
        and adds correctness hazards (unbounded gather, blocking calls in the
        loop).
\end{enumerate}

\subsection{Staff}

\begin{enumerate}
  \setcounter{enumi}{30}
  \item \textbf{Q: Design an organization-wide HTTP client policy. What is in
        it and how is it enforced?}\\
        \textit{A:} Mandated: four-dimension timeouts with ceilings, jittered
        bounded retries plus a shared retry budget, idempotency gates on
        mutation, header-safe telemetry, hermetic MockTransport tests.
        Enforced as a shared internal client library (policy as code), CI
        lint rules banning bare timeout-free calls, and fault-injection tests
        in staging. Policy documents that CI cannot check are fiction.

  \item \textbf{Q: A downstream provider publishes a 100 req/s limit. Your
        fleet of 200 pods each has a local token bucket set to 1 req/s.
        What goes wrong and how do you architect past it?}\\
        \textit{A:} Local buckets under-utilize (fairness is emergent, not
        guaranteed) and can't react to provider-side repartitioning; pod
        autoscaling silently raises aggregate rate. Architecture: make local
        limits a function of live fleet size (or a central rate-limiting
        sidecar), treat \texttt{429}s as control signals that back the whole
        fleet off, and negotiate limits as an SLO with the provider rather
        than reverse-engineering them from error rates.
\end{enumerate}

%% ---------------------------------------------------------------------------
%% 9. External Bibliography & Further Reading
%% ---------------------------------------------------------------------------

\section{External Bibliography \& Further Reading}
\label{sec:ch02-bibliography}

\begin{itemize}
  \item \textbf{httpx documentation} \cite{httpxdocs} --- the authoritative
        reference for the four-dimension timeout model, \texttt{Limits}
        pooling, event hooks, HTTP/2 opt-in, and \texttt{MockTransport}.
        Read the ``Timeouts'' and ``Advanced'' pages in full.
  \item \textbf{Requests documentation} \cite{requestsdocs} --- the Session
        abstraction and the historical baseline every Python engineer
        inherits; its ``Advanced Usage'' chapter remains excellent on
        connection pooling semantics.
  \item \textbf{RFC 9110, \emph{HTTP Semantics}} \cite{rfc9110} --- the
        normative definitions of method idempotency (\S9.2.2), status-code
        classes, \texttt{429}/\texttt{503}, and \texttt{Retry-After}
        (\S15.5.4, \S10.2.3). The primary source for every retry rule in
        this chapter.
  \item \textbf{RFC 8288, \emph{Web Linking}} --- the \texttt{Link} header
        and \texttt{rel="next"} consumed in
        Listing~\ref{lst:ch02-pagination}.
  \item \textbf{Michael T. Nygard, \emph{Release It!}, 2nd ed.}
        \cite{nygard2018releaseit} --- the canonical treatment of cascading
        failure, timeouts, and circuit breakers. Chapter 4 (``Stability
        Antipatterns'') is the war story's playbook.
  \item \textbf{Marc Brooker, AWS Architecture Blog, \emph{``Exponential
        Backoff and Jitter''} and \emph{``Timeouts, retries, and backoff
        with jitter''}} --- the industry-standard derivations of full jitter
        and decorrelated jitter that Section~\ref{sec:ch02-retries}
        formalizes. \url{https://aws.amazon.com/blogs/architecture/exponential-backoff-and-jitter/}
  \item \textbf{Tenacity documentation} --- the retry library used in
        Listing~\ref{lst:ch02-retry}: composable stop/wait/retry strategies.
        \url{https://tenacity.readthedocs.io/}
  \item \textbf{pydantic documentation} \cite{pydanticdocs} --- typed
        response models and validation for the SDK-wrapper pattern; Chapter~4
        develops contract-first serialization in depth.
  \item \textbf{The Twelve-Factor App} \cite{twelvefactor} --- factor III
        (config in the environment) underpins the credential-handling rules
        of Section~\ref{sec:ch02-pitfalls}.
  \item \textbf{Martin Kleppmann, \emph{Designing Data-Intensive
        Applications}} \cite{kleppmann2017ddia} --- Chapter 8's treatment of
        partial failure and ``we cannot know'' outcomes at network
        boundaries; the theoretical frame for our error taxonomy.
\end{itemize}

%% ---------------------------------------------------------------------------
%% 10. Capstone Projects
%% ---------------------------------------------------------------------------

\section{Capstone Projects}
\label{sec:ch02-capstones}

Five progressive projects. All run offline against the Northwind stub
(Listing~\ref{lst:ch02-stub}) or \texttt{httpx.MockTransport}; all require a
passing pytest suite; all follow the workspace conventions (typed Python
3.11+, PowerShell docs, synthetic data only).

\subsection{Capstone 2.1 [junior] --- Robust CLI Downloader}

\textbf{Objective.} Build \texttt{nrget}: a PowerShell-friendly CLI that
downloads a part's datasheet (\texttt{/catalog/export}) from the stub to a
local file, correctly and observably.

\textbf{Inputs/Datasets.} The local stub's \texttt{/catalog/export} endpoint
(synthetic Northwind catalog, NDJSON).

\textbf{Steps.}
\begin{enumerate}
  \item Wrap \texttt{httpx.Client} with explicit four-dimension timeouts and
        a typed error taxonomy (reuse Listing~\ref{lst:ch02-httpx} patterns).
  \item Stream the body with \texttt{iter\_bytes}, a \texttt{.part} temp
        file, and atomic rename on success.
  \item Print byte-level progress to the console; support a
        \texttt{--max-bytes} guard.
  \item Exit codes: 0 success, 2 client error, 3 server/transient error
        after bounded retries.
\end{enumerate}

\textbf{Acceptance Criteria.}
\begin{itemize}
  \item[$\square$] No call lacks an explicit timeout (verified by code
        inspection and a test asserting the client's \texttt{timeout} fields).
  \item[$\square$] A killed download leaves no partial artifact.
  \item[$\square$] MockTransport tests cover 200, 404, 500-then-200, and
        malformed-body paths; suite passes with the stub stopped.
  \item[$\square$] Logs contain method/path/status/latency and no headers.
\end{itemize}

\textbf{Stretch Goals.} Resume support via \texttt{Range} requests; a
\texttt{--dry-run} mode that prints the request plan without sending.

\subsection{Capstone 2.2 [mid] --- Rate-Limit-Aware Catalog Crawler}

\textbf{Objective.} Crawl the entire paginated catalog of an instrumented
stub that enforces a published limit of 5 req/s and answers
\texttt{429 + Retry-After} when exceeded --- without ever being throttled in
steady state.

\textbf{Inputs/Datasets.} Extended stub: cursor-paginated \texttt{/parts}
over a 500-part synthetic catalog, a token-bucket limiter, and
\texttt{RateLimit-Policy}/\texttt{Retry-After} headers.

\textbf{Steps.}
\begin{enumerate}
  \item Implement a client-side token bucket ($r = 4$ tokens/s, $B = 4$)
        gating every request.
  \item Wrap calls in the retry policy of Listing~\ref{lst:ch02-retry};
        \texttt{Retry-After} overrides jitter.
  \item Consume pages through the cursor iterator of
        Listing~\ref{lst:ch02-pagination}.
  \item Record per-request timestamps to a CSV; compute achieved rate.
\end{enumerate}

\textbf{Acceptance Criteria.}
\begin{itemize}
  \item[$\square$] Full crawl completes; zero \texttt{429}s in the steady
        state (allow one cold-start burst violation, documented).
  \item[$\square$] Achieved sustained rate $\le$ 5 req/s, evidenced from the
        timestamp CSV (not wall-clock test assertions).
  \item[$\square$] A forced-\texttt{429} MockTransport test proves
        \texttt{Retry-After} overrides jitter exactly.
  \item[$\square$] Deterministic suite: seeded RNG, injected sleep, no
        network.
\end{itemize}

\textbf{Stretch Goals.} Adaptive bucket that lowers $r$ on 429 and recovers
slowly (additive-increase); parallel crawlers sharing one process-wide
bucket.

\subsection{Capstone 2.3 [mid] --- Typed SDK Wrapper with Full Test Matrix}

\textbf{Objective.} Publish (locally) a \texttt{northwind-parts-sdk} package:
typed \texttt{pydantic} response models, the resilient client, pagination
iterators, and a complete offline fault matrix.

\textbf{Inputs/Datasets.} The Northwind stub endpoints; a hand-written JSON
Schema for the \texttt{Part} resource used to generate/validate the
\texttt{pydantic} model.

\textbf{Steps.}
\begin{enumerate}
  \item Model \texttt{Part}, \texttt{PartPage}, and \texttt{ApiError} as
        \texttt{pydantic} models with strict types \cite{pydanticdocs}.
  \item Parse responses into models; map validation failures to
        \texttt{ApiContractError} with field-level context.
  \item Compose timeouts, pooling, auth injection, hooks, and the tenacity
        retry wrapper behind a clean public API.
  \item Build the MockTransport matrix: every status class, transport
        timeouts, mid-stream reset, truncated JSON, wrong
        \texttt{Content-Type}, oversized page count.
\end{enumerate}

\textbf{Acceptance Criteria.}
\begin{itemize}
  \item[$\square$] Public API is fully typed; \texttt{mypy}-clean (or
        equivalent checked annotations) on the package.
  \item[$\square$] Fault matrix covers $\ge$ 15 scenarios including every
        taxonomy branch; all deterministic and offline.
  \item[$\square$] Contract errors are never retried; transient errors retry
        within budget; idempotency gate enforced on mutations.
  \item[$\square$] README (PowerShell-only instructions) documents install,
        usage, and error semantics.
\end{itemize}

\textbf{Stretch Goals.} Generate the model from the JSON Schema at build
time; add an async facade reusing the same core via
\texttt{httpx.AsyncClient}.

\subsection{Capstone 2.4 [senior] --- Chaos-Client Harness}

\textbf{Objective.} Build a fault-injection transport that sits between your
resilient client and the stub and proves, empirically, that the resilience
policy works: the client survives a gauntlet of injected faults with
predictable behavior.

\textbf{Inputs/Datasets.} A \texttt{ChaosTransport} wrapping
\texttt{httpx.MockTransport} with a scripted fault schedule: connect resets,
read stalls, truncation, 500 bursts, 429s with/without \texttt{Retry-After},
gigantic \texttt{Content-Length} lies, and gzip bombs (synthetic, capped).

\textbf{Steps.}
\begin{enumerate}
  \item Implement \texttt{ChaosTransport(BaseTransport)} reading a YAML/JSON
        fault schedule; deterministic via seeded RNG.
  \item Define the expected behavior per fault (recover, fail-fast with
        typed error, abort with bounded time) as an executable contract.
  \item Run the full SDK against each schedule; assert both outcome
        \emph{and} bounds: total attempts, total sleep, bytes buffered.
  \item Produce a coverage report mapping faults $\rightarrow$ observed
        client behavior.
\end{enumerate}

\textbf{Acceptance Criteria.}
\begin{itemize}
  \item[$\square$] $\ge$ 8 distinct fault classes, each with an executable
        expectation.
  \item[$\square$] Every run is reproducible from its seed; no wall-clock
        dependence.
  \item[$\square$] At least one fault deliberately exceeds the retry budget
        and the client demonstrably fails fast with the correct typed error.
  \item[$\square$] Decompression-bomb test proves the streaming path's size
        guard engages.
\end{itemize}

\textbf{Stretch Goals.} Latency distributions (injected p99 spikes) with a
fake clock; chaos \emph{schedules} replayed from recorded production incident
timelines (synthetic).

\subsection{Capstone 2.5 [staff] --- Async Concurrent Fetcher with Bounded Concurrency}

\textbf{Objective.} Design and validate an \texttt{httpx.AsyncClient}-based
fetcher that drains a 10{,}000-item work queue from the stub at maximum
throughput \emph{without} tripping the provider's rate limit, exhausting the
pool, or losing backpressure --- and write the design memo defending every
knob.

\textbf{Inputs/Datasets.} Instrumented stub with a 200-part synthetic
catalog, per-IP token bucket at 50 req/s, and injected latency jitter
(5--200~ms, seeded).

\textbf{Steps.}
\begin{enumerate}
  \item Bounded concurrency via \texttt{asyncio.Semaphore} (not unbounded
        \texttt{gather}); pool limits aligned to the semaphore.
  \item Process-wide token bucket below the published limit; \texttt{429}
        handling that backs off the \emph{whole} fetcher, not one task.
  \item Retry policy per Listing~\ref{lst:ch02-retry} semantics, adapted for
        async sleep; elapsed and attempt budgets enforced.
  \item Cancellation: a single Ctrl+C (or task cancel) drains in-flight
        work, closes streams, and exits cleanly with a partial-results
        manifest.
  \item The memo: Little's-law sizing of concurrency vs.\ rate limit; pool
        vs.\ semaphore alignment; failure-mode table.
\end{enumerate}

\textbf{Acceptance Criteria.}
\begin{itemize}
  \item[$\square$] Full queue drains with zero unhandled exceptions and zero
        pool timeouts at the chosen settings.
  \item[$\square$] Sustained rate $\le$ 50 req/s evidenced from stub-side
        logs; 429 rate $<$ 0.1\% after warm-up.
  \item[$\square$] Deterministic tests via \texttt{MockTransport} + injected
        async sleep; no wall-clock assertions.
  \item[$\square$] Cancellation test proves no orphan tasks and no partial
        artifacts.
  \item[$\square$] Design memo quantifies every knob; trade-offs and
        rejected alternatives documented.
\end{itemize}

\textbf{Stretch Goals.} Adaptive concurrency controller (AIMD on 429/latency
signals); fairness across multiple upstream hosts with per-host buckets.

%% ---------------------------------------------------------------------------
%% 11. Python Dependencies Appendix
%% ---------------------------------------------------------------------------

\section{Python Dependencies Appendix}
\label{sec:ch02-dependencies}

The chapter's code pins four third-party packages. Python 3.11+ is required
throughout. Install on Windows/PowerShell:

\begin{minted}{bash}
python -m venv .venv
.\.venv\Scripts\Activate.ps1
pip install -r requirements.txt
\end{minted}

\begin{minted}[caption={\texttt{requirements.txt} for Chapter 2.},label={lst:ch02-requirements}]{text}
httpx>=0.27,<1.0
tenacity>=8.3
pydantic>=2.7,<3.0
pytest>=8.0
\end{minted}

\subsection{httpx ($\geq$ 0.27)}

\textbf{Description.} A fully typed HTTP client offering a
\texttt{requests}-compatible API with synchronous and asynchronous clients,
optional HTTP/2, four-dimension timeouts, connection pooling via
\texttt{Limits}, event hooks, and \texttt{MockTransport} for hermetic tests
\cite{httpxdocs}.

\textbf{Why included.} It is the chapter's foundation: the only mainstream
Python client that unifies sync/async, expresses the full timeout taxonomy of
Section~\ref{sec:ch02-timeouts}, and makes offline testing a first-class
feature rather than a monkey-patch.

\textbf{Alternatives.} \texttt{requests} \cite{requestsdocs} (sync-only,
HTTP/1.1-only, two-dimension timeouts); \texttt{aiohttp} (async-only, its own
API shape); stdlib \texttt{urllib} (dependency-free but austere, see
Section~\ref{sec:ch02-evolution}).

\textbf{Installation.} \texttt{pip install "httpx>=0.27,<1.0"} --- pin the
upper bound while the library remains pre-1.0.

\subsection{tenacity ($\geq$ 8.3)}

\textbf{Description.} A general-purpose retrying library: composable
\texttt{stop}, \texttt{wait}, and \texttt{retry} strategies, with support for
coroutines, statistics, and custom sleep injection.

\textbf{Why included.} Correct retry loops are subtle (budget arithmetic,
exception filtering, jitter distribution, sleep abstraction). Tenacity makes
the policy declarative and --- decisive for this book --- accepts an
injected \texttt{sleep}, enabling the deterministic fault-matrix tests of
Section~\ref{sec:ch02-tests}.

\textbf{Alternatives.} \texttt{backoff} (decorator-oriented, less
composable); \texttt{urllib3.util.Retry} (only within requests/urllib3, and
its \texttt{respect\_retry\_after\_header} defaults helpfully apply there);
hand-rolled loops (acceptable for one-off scripts, error-prone at scale).

\textbf{Installation.} \texttt{pip install "tenacity>=8.3"}.

\subsection{pydantic ($\geq$ 2.7)}

\textbf{Description.} Data-validation library using Python type annotations:
declares typed models, validates and coerces inbound data, and produces
structured errors with field-level paths \cite{pydanticdocs}.

\textbf{Why included.} The SDK-wrapper pattern (Capstone 2.3) needs a
contract layer that turns ``200 OK with a surprising body'' into a precise,
actionable \texttt{ApiContractError}. pydantic v2's Rust core makes
validation cheap enough for the hot path.

\textbf{Alternatives.} \texttt{dataclasses} + manual checks (stdlib, no
validation); \texttt{msgspec} (faster, smaller ecosystem); \texttt{attrs} +
\texttt{cattrs} (mature, more assembly required).

\textbf{Installation.} \texttt{pip install "pydantic>=2.7,<3.0"}.

\subsection{pytest ($\geq$ 8.0)}

\textbf{Description.} The standard Python test framework: fixtures,
parameterization, rich assertion introspection, and a plugin ecosystem.

\textbf{Why included.} Every listing in this chapter is validated by the
fault-matrix suite (Listing~\ref{lst:ch02-tests}); pytest's
\texttt{raises} introspection and fixtures make exception-taxonomy tests
terse and readable.

\textbf{Alternatives.} \texttt{unittest} (stdlib, more boilerplate);
\texttt{nose2} (smaller ecosystem). For async tests add
\texttt{pytest-asyncio} --- deliberately excluded here to keep the chapter's
suite purely synchronous and deterministic.

\textbf{Installation.} \texttt{pip install "pytest>=8.0"}.

%% ---------------------------------------------------------------------------
%% Chapter Summary & Key Takeaways
%% ---------------------------------------------------------------------------

\section{Chapter Summary \& Key Takeaways}

\begin{itemize}
  \item \textbf{Client choice is architectural.} \texttt{urllib} for
        zero-dependency bootstrap code; \texttt{requests} for legacy
        synchronous HTTP/1.1 estates; \texttt{httpx} for everything new ---
        async, HTTP/2, four-dimension timeouts, and MockTransport testability
        \cite{httpxdocs,requestsdocs}.
  \item \textbf{Every wait is bounded or it is a bug.} Connect, read, write,
        and pool timeouts bound four distinct hangs; ``no timeout'' is the
        enabling defect of cascading failure \cite{nygard2018releaseit}.
  \item \textbf{Pools are capacity.} One long-lived client per base URL;
        size \texttt{max\_connections} to peak concurrency and
        \texttt{max\_keepalive\_connections} to steady state; churn re-pays
        DNS/TCP/TLS per request.
  \item \textbf{Retries are load.} Cap attempts \emph{and} elapsed time,
        jitter fully (expected total wait $\frac{b}{2}(2^N - 1)$), honor
        \texttt{Retry-After}, budget the fleet, and gate non-idempotent
        methods behind idempotency keys \cite{rfc9110}.
  \item \textbf{Consume pagination lazily and payloads as streams.}
        Iterators hide offset/cursor/link-header mechanics;
        \texttt{iter\_bytes} keeps memory $O(\text{chunk})$; cancellation
        and atomic rename keep artifacts clean.
  \item \textbf{Classify before you react.} Transport, status, and contract
        failures have different owners and remedies; parse only after the
        status check; never retry deterministic failures.
  \item \textbf{Test hermetically.} \texttt{MockTransport}, seeded
        randomness, and injected sleep yield a deterministic, offline fault
        matrix --- the only kind of test suite a consumer library deserves.
\end{itemize}

%% ---------------------------------------------------------------------------
%% Quick-Reference Card
%% ---------------------------------------------------------------------------

\section{Quick-Reference Card}
\label{sec:ch02-quickref}

\subsection{Timeout Cheat-Sheet}

\begin{table}[h]
  \centering
  \small
  \begin{tabularx}{\textwidth}{l l X l}
    \toprule
    \textbf{Timeout} & \textbf{Bounds} & \textbf{Failure it catches} & \textbf{Starting value} \\
    \midrule
    connect & DNS + TCP + TLS setup & dead/firewalled host, TLS misconfig & 2~s \\
    read    & waiting for response bytes & hung app behind healthy LB; stalled stream & 10~s (align to endpoint p99) \\
    write   & sending request body & peer that stopped reading (uploads) & 10~s \\
    pool    & waiting for a free connection & concurrency $>$ pool size (queueing) & 2~s \\
    \bottomrule
  \end{tabularx}
  \caption{Timeout dimensions: always set all four; never leave any infinite.}
  \label{tab:ch02-timeout-cheatsheet}
\end{table}

\subsection{Retry Decision Table}

\begin{table}[h]
  \centering
  \small
  \begin{tabularx}{\textwidth}{l l X}
    \toprule
    \textbf{Signal} & \textbf{Retry?} & \textbf{Conditions} \\
    \midrule
    Connect/Read timeout, connection reset & Yes & idempotent method or idempotency key; jittered backoff \\
    408, 425 & Yes & same as above \\
    429 & Yes & after \texttt{Retry-After} (clamped); back off the bucket \\
    500, 502, 503, 504 & Yes & bounded attempts + elapsed budget; 503 may carry \texttt{Retry-After} \cite{rfc9110} \\
    400, 401, 403, 404, 409, 422 & No & deterministic: fix the request, rotate credentials, or reconcile state \\
    2xx with unparseable/invalid body & No & \texttt{ApiContractError}: provider bug or version skew --- alert \\
    POST/PUT-side-effect without idempotency key & No & replay may duplicate effects; escalate to caller \\
    \bottomrule
  \end{tabularx}
  \caption{Retry decisions by failure signal.}
  \label{tab:ch02-retry-decisions}
\end{table}

\subsection{Status-Class Handling Table}

\begin{table}[h]
  \centering
  \small
  \begin{tabularx}{\textwidth}{l l X}
    \toprule
    \textbf{Class} & \textbf{Meaning} & \textbf{Consumer action} \\
    \midrule
    1xx & Interim & handled by the transport; never surfaced to callers \\
    2xx & Success & then (and only then) parse; validate against the contract \\
    3xx & Redirect & follow per client policy; preserve method semantics; cap hops \\
    4xx & Caller fault & typed \texttt{ApiClientError4xx}; fix request; never retry as-is \\
    5xx & Server fault & typed \texttt{ApiServerError5xx}; retry per policy; alert on sustained rate \\
    \bottomrule
  \end{tabularx}
  \caption{Status-class handling for consumers (RFC 9110 \cite{rfc9110}).}
  \label{tab:ch02-status-classes}
\end{table}

\section{Standardized Evidence Addendum}
\subsection{Benchmark Snapshot Template}
\begin{table}[h]
  \centering
  \small
  \begin{tabularx}{\textwidth}{l c c c c}
    \toprule
    \textbf{Config} & \textbf{p50 latency} & \textbf{p99 latency} & \textbf{Error rate} & \textbf{Throughput} \\
    \midrule
    Baseline & -- & -- & -- & 1.00x \\
    Candidate A & -- & -- & -- & -- \\
    Candidate B & -- & -- & -- & -- \\
    \bottomrule
  \end{tabularx}
  \caption{Minimum benchmark table to keep per chapter.}
\end{table}

\subsection{Request Trace Template}
\begin{itemize}
  \item \textbf{Request:} one representative API call for this chapter's topic.
  \item \textbf{Before:} behavior (headers, payload, latency, failure mode) with the naive approach.
  \item \textbf{After:} behavior with the technique in this chapter.
  \item \textbf{Result:} what changed in correctness, resilience, latency, and operability.
\end{itemize}

\subsection{Cost and Latency Budget Snapshot}
\begin{table}[h]
  \centering
  \small
  \begin{tabularx}{\textwidth}{l c c}
    \toprule
    \textbf{Stage} & \textbf{Latency budget} & \textbf{Cost driver} \\
    \midrule
    TLS / connection setup & -- & handshake round trips, keep-alive hit rate \\
    Request validation & -- & schema complexity, CPU per request \\
    Business logic and I/O & -- & downstream fan-out, query cost \\
    Serialization and egress & -- & payload size, compression ratio \\
    \bottomrule
  \end{tabularx}
  \caption{Operational budget template for this chapter's technique.}
\end{table}

\section{Configuration Preset Template}
\begin{minted}{text}
# api_policy.yaml
policy_version: "v1.0.0"
component: "<chapter_topic>"
timeout_ms: 0
retry_budget: 0
fallback_strategy: "fail_fast"
rollback_trigger: "error_budget_burn_gt_threshold"
\end{minted}

\section{CI Quality Gate Specification}
\begin{itemize}
  \item contract tests green against the pinned OpenAPI/AsyncAPI/Protobuf spec,
  \item no breaking schema changes without a version bump,
  \item error responses conform to RFC 9457 problem-details shape,
  \item benchmark deltas within approved regression bounds (latency and throughput),
  \item policy version and rollback plan present in release metadata.
\end{itemize}

\section{Incident Runbook Template}
\begin{enumerate}
  \item Detect regression from SLO burn-rate alerts or consumer reports.
  \item Scope impacted endpoints, versions, and consumer segments.
  \item Contain impact (rollback, feature flag, rate-limit tightening, or traffic shift).
  \item Diagnose with traces, structured logs, and contract diffs.
  \item Recover with validated fix and verified rollout procedure.
  \item Prevent recurrence via new regression tests and CI gates.
\end{enumerate}

\section{Cross-Chapter Decision Card}
\begin{table}[h]
  \centering
  \small
  \begin{tabularx}{\textwidth}{l X X}
    \toprule
    \textbf{Dimension} & \textbf{Choose this approach when} & \textbf{Prefer an alternative when} \\
    \midrule
    Use case & -- & -- \\
    Latency budget & -- & -- \\
    Operational cost & -- & -- \\
    Team maturity & -- & -- \\
    \bottomrule
  \end{tabularx}
  \caption{Decision card to complete per chapter.}
\end{table}

\section{ROI Model and Build-vs-Buy}
\begin{itemize}
  \item \textbf{Build when:} the capability is differentiating, compliance forbids vendors, or scale makes vendor pricing irrational.
  \item \textbf{Buy when:} the capability is commodity (auth, gateways, observability), time-to-market dominates, or staffing is thin.
  \item \textbf{Hidden costs to model:} on-call load, upgrade cadence, expertise ramp, data egress fees.
\end{itemize}

\section{Disaster Recovery Notes (RPO/RTO)}
\begin{itemize}
  \item Define recovery point objective (RPO): maximum acceptable data loss window.
  \item Define recovery time objective (RTO): maximum acceptable restoration time.
  \item Test restores regularly; an untested backup is not a backup.
\end{itemize}


